% =========================================================================
%  Toward Relational Being in Intimate Life:
%  Integrative Justice and an Ethics of Dialectical-Positivist Practice
%  — A Family-and-Relationship Management Application as a Case Study
%
%  Wanhong Huang (Independent researcher)
%
%  Compile with XeLaTeX:
%      xelatex main.tex && biber main && xelatex main.tex && xelatex main.tex
%  or:
%      latexmk -xelatex main.tex
% =========================================================================

\documentclass[11pt,a4paper]{article}
\usepackage{serendip-paper}

% ── Bibliography (add refs.bib to this folder) ───────────────────────
\usepackage[backend=biber, style=authoryear, sorting=nyt, maxcitenames=2]{biblatex}
\addbibresource{../../bib/master.bib}

\newcommand{\TODO}[1]{\textcolor{gold}{\textsf{[TODO: #1]}}}

% ── Title page is set manually after \begin{document} ────────────────

\begin{document}

% ── Title page (centered) ────────────────────────────────────────────
\thispagestyle{empty}
\null\vfill
\begin{center}
{\footnotesize\color{rosedeep}\scshape Philosophy of Intimacy and the Theory of Justice \quad$\cdot$\quad Paper I\par}
\vspace{0.9em}
\coverrule\par
\vspace{1.6em}
{\color{warnred}\Large Toward Relational Being in Intimate Life: Integrative Justice and an Ethics of Dialectical-Positivist Practice\par}
\vspace{1.0em}
{\color{warnred}\large\itshape A Family-and-Relationship Management Application as a Case Study\par}
\vspace{2.4em}
\coverrule[0.25]\par
\vspace{1.4em}
{\large Wanhong Huang\par}
\vspace{0.5em}
{\footnotesize \texttt{huangwanhong@serendip.ngo}\par}
\vspace{1.2em}
{\today\par}
\end{center}
\vfill\null
\coverfootnote
\clearpage

% ── Dedication ───────────────────────────────────────────────────────
\clearpage
\thispagestyle{empty}
\vspace*{0.20\textheight}
\begin{center}
{\itshape\color{ink}
To the one I love deeply,\\[0.4em]
a fellow native of the homeland we both left behind,\\[0.25em]
yet whom I was to meet, by some miracle, for the first time\\[0.25em]
only in a country far from the one that bore us.\\}
\vspace{2.6em}
{\cjkfont\large\color{rosedeep} 两情若是久长时，又岂在朝朝暮暮}\\[0.8em]
{\itshape\small\color{rosedeep}
If two hearts are to hold each other for a lifetime,\\[0.18em]
why should they grieve to be apart\\[0.18em]
from one morning, one evening, to the next?}\\[1.0em]
{\footnotesize\color{lightgold}after Qin Guan (\cjkfont 秦观\/{\normalfont,} 1049--1100)}
\vspace{2.4em}

{\cjkfont\large\color{rosedeep} 但愿人长久，千里共婵娟}\\[0.8em]
{\itshape\small\color{rosedeep}
May we both be granted long years,\\[0.18em]
and though the whole world lie between us,\\[0.18em]
may we look up, each from afar,\\[0.18em]
upon the same bright moon.}\\[1.0em]
{\footnotesize\color{lightgold}after Su Shi (\cjkfont 苏轼\/{\normalfont,} 1037--1101)}
\end{center}
\vfill
\clearpage

% =========================================================================
\begin{abstract}
\noindent
Positivist, evidence-based technologies increasingly mediate intimate life: couples coordinate, record, and remember through software whose default logic is measurement and optimization. This paper asks how such technologies ought instead to be governed if their aim is the flourishing of a relationship, the sustaining of what it calls \emph{relational being} in intimate partnership.

\smallskip
\noindent It makes four kinds of claim. (1) \emph{Normative}: an \emph{integrative-justice} framework, resting on a relational ontology, draws on a dozen ethical and epistemic traditions and assigns each the questions it is competent to answer, in a defended ordering rather than a numerical average. (2) \emph{Practical}: an \emph{ethics of dialectical-positivist practice} warrants measurement differentially, site by site, according to whether a phenomenon is constituted partly by remaining unmeasured. (3) \emph{Engineering}: a working application, designed by the author, instantiates the framework in software, with event-driven household coordination, a witness-only AI layer, co-creative exploration, and the deliberate \emph{absence} of AI where it would intrude. (4) \emph{Diagnostic}: it names wrongs specific to intimate computing, among them \emph{proxy epistemic injustice}, the misattribution of an act of knowing, and the opposition between structural efficiency and relational development.

\smallskip
\noindent The limits are explicit. (i) Mediation's deepest harms, relational alienation and epistemic distortion, can be reduced but never wholly removed, so the framework offers harm-reduction and a discipline of restraint. (ii) Its dialectical adjudication is a practice of judgement for which no algorithm is given. (iii) The evidence is a single, author-designed case, offered as illustration and existence proof. The framework is meant, finally, to bear on adjacent settings as well: the ethics of the societal use of data, and the design of mechanisms linking the family to public institutions.

\vspace{0.6em}
\noindent\textbf{Keywords}\quad
relational being; intimate computing; integrative justice; the ethics of dialectical-positivist practice; generative presence; data ethics; the ethics of flourishing.
\end{abstract}

% =========================================================================
\section{Introduction}\label{p01:sec:intro}

A growing share of intimate life is now conducted through, or recorded by, software. Couples coordinate on shared calendars, keep joint ledgers, log meals and moods, and increasingly delegate to conversational agents the small administrative friction of a shared existence. The literature that examines this development has been overwhelmingly either empirical, concerned with how couples in fact use such tools, or design-methodological, concerned with how such tools should be built. The normative question, how \emph{ought} the positivist, evidence-based mediation of intimate life be governed and toward what end, has received comparatively little sustained philosophical treatment. This paper is an attempt at that treatment.

The end proposed here is not efficiency, nor wellbeing as conventionally operationalized, but \emph{relational flourishing}: the sustaining over time of a relationship understood, after the relational-ontological tradition, as ontologically prior to the individuals it relates. The paper restricts this large idea deliberately. It does not advance a general theory of relational being; it asks only what relational being requires in the specific domain of \emph{intimate partnership}, and what follows for the technologies admitted into that domain. The restriction matters because intimate partnership has features that distinguish it from the family at large, the workplace, and the polity, and that bear directly on which forms of measurement and mediation are admissible.

The paper's central constructive proposal is an \emph{integrative-justice} framework. By this is meant a way of bringing several normative traditions (evidentialism in epistemology, the positivism--anti-positivism dialectic in the philosophy of the human sciences, deontological ethics, feminist care ethics, Confucian role-ethics and Daoist non-interference, and the capability approach to flourishing) to bear jointly on the ethics of intimate data, in such a way that each tradition is assigned the questions it can competently answer and denied the questions it cannot. The integration is not eclectic averaging; it is a division of normative labour. Evidentialism governs what the system is entitled to claim to know; deontology fixes the inviolable constraints on the treatment of the other; care ethics governs the texture of responsiveness; the capability approach supplies the criterion of flourishing against which the whole is measured; and the East Asian sources supply an account, largely absent from the Western canon, of the positive value of restraint and non-action.

From this framework the paper derives a practical posture it calls an \emph{ethics of dialectical-positivist practice}. The positivist holds that domestic phenomena are legitimately quantifiable and therefore improvable through measurement. The anti-positivist, drawing on Diltheyan, Weberian, Geertzian, and Frankfurt-school sources, denies that the constitutive phenomena of intimate life submit without distortion to such operationalization. The posture defended here treats this opposition dialectically: measurement is neither universally licensed nor universally corrosive, but warranted differentially, site by site, according to whether the phenomenon at issue is partly constituted by \emph{not} being measured. The dialectic is not resolved once and for all; it is re-adjudicated at each site of possible measurement, and the adjudication itself is treated as a first-order ethical act.

To keep the argument from floating free of practice, the paper takes as a sustained case study a working family-and-relationship management application that the author designed and maintains: roughly two dozen interlocking modules, among them a shared journal, a household ledger, a health log, a synchronized calendar, a conversational AI layer, and an etiquette aid, each of which embodies, in concrete software, a determinate answer to the question of where on the dialectical continuum it should sit. The application is used here not as the object of a design-research report but as an extended worked example: a place where the framework's abstract commitments can be seen taking on operational form, and where their costs and tensions become visible.

One feature of the case is methodologically and ethically unusual, and the paper takes it up as a philosophical instrument. The application has been built for an intimate relationship in which the other party is, at the time of writing, unaware of the application's existence. Designing for an absent or unknowing other is a limiting case that sharpens every question a relational ethics must answer: What may one party legitimately decide on behalf of a relationship the other has not been consulted about? When does preparation for another's good become presumption upon it? What is owed to a person who is the subject of a system they have not consented to? The paper argues that the integrative-justice framework yields determinate and sometimes uncomfortable answers to these questions, answers that constrain the project significantly, and that a framework which could not address the absent-other case would be inadequate to the ordinary case as well.

The paper is accordingly organized around three questions. The first concerns normative foundations: what does relational flourishing, in intimate partnership specifically, require of the technologies admitted into that relationship, and which normative traditions are competent to specify those requirements? The second concerns integration: how may evidentialist, positivist and anti-positivist, deontological, care-ethical, Confucian-Daoist, and capability-theoretic considerations be integrated into a single coherent practice rather than left as competing and incommensurable demands? The third concerns the limiting case of the absent other: what does this framework imply for the case of designing for an intimate other who is unaware of, and has not consented to, the design, and what does that limiting case reveal about the ordinary one?

The paper's contributions are correspondingly philosophical rather than technical, and may be stated as follows: (i) a restriction and specification of the relational-being thesis to the domain of intimate partnership (§\ref{p01:sec:relational-being}); (ii) an integrative-justice framework for the ethics of intimate data, with each constituent tradition given an explicit competence and an explicit limit (§\ref{p01:sec:frameworks}--§\ref{p01:sec:integration-argument}); (iii) the articulation of an ethics of dialectical-positivist practice as the framework's practical upshot (§\ref{p01:sec:dialectic}); (iv) a worked case study showing the framework instantiated in a real system (§\ref{p01:sec:casestudy}); (v) the identification of wrongs specific to intimate computing, including \emph{proxy epistemic injustice}, the misattribution of an act of knowing, and the opposition between structural efficiency and relational development (§\ref{p01:sec:calendar}, §\ref{p01:sec:limitations}); (vi) a treatment of the absent-other problem as a test of relational ethics in general (§\ref{p01:sec:absent-other}); and (vii) a set of implications reaching beyond intimate partnership, for the ethics of the societal use of data and the design of mechanisms linking the family to public institutions (§\ref{p01:sec:conclusion}). Questions of software engineering, such as architecture, storage, and choice of frameworks, are treated only insofar as they bear on the normative argument, and are confined to a brief appendix; they are not the subject of this paper.

The remainder proceeds as follows. §\ref{p01:sec:relational-being} specifies relational being for intimate partnership and situates its intellectual lineage. §\ref{p01:sec:frameworks} sets out the constituent normative traditions, one per subsection. §\ref{p01:sec:integration-argument} argues for their integration. §\ref{p01:sec:dialectic} develops the dialectical-positivist posture. §\ref{p01:sec:casestudy} presents the case study. §\ref{p01:sec:absent-other} takes up the absent other. §\ref{p01:sec:limitations} states limits. §\ref{p01:sec:related} places the paper within the surrounding literatures, and §\ref{p01:sec:conclusion} concludes.

% =========================================================================
% =========================================================================
% =========================================================================
\section{Relational Being in Intimate Partnership}\label{p01:sec:relational-being}

\subsection{Intellectual Lineage}\label{p01:sec:lineage}

The ontological commitments developed in this section are post-phenomenological in the lineage of \textcite{verbeek2005things}, who reads artefacts as mediating rather than merely containing the practices around them, and relational in the lineages of \textcite{buber1937iandthou}, for whom the \emph{I-Thou} relation is ontologically primary, \textcite{ames2011confucian} and \textcite{tu1985confucian}, for whom the self is constituted in its roles and relations, and \textcite{borgmann1984technology}, who develops the concept of \emph{focal practices} whose value is intrinsic and erodable by instrumentalization. What this paper adds is a restriction of these large ontological claims to the specific domain of intimate partnership, and an argument that the restriction changes what the ontology demands. A fuller placement of the paper within the literatures of intimate technology, self-quantification, and autobiographical design is deferred to §\ref{p01:sec:related}, after the framework it positions has been set out.

\subsection{The Ontological Ground: Relationality and Generative Presence}\label{p01:sec:ontology}

The framework rests on a relational ontology, which the present author has developed at length elsewhere \parencite{huang2025relational} and which converges with a wide tradition: Buber's primacy of the \emph{I-Thou} relation \parencite{buber1937iandthou}, Gergen's relational being \parencite{gergen2009relational}, Whitehead's process metaphysics of relational ``occasions of experience'' \parencite{whitehead1978process}, Heidegger's \emph{Mitsein} or being-with \parencite{heidegger1962being}, the Buddhist doctrine of dependent origination and emptiness \parencite{nagarjuna1995mulamadhyamakakarika}, and Confucian role-ethics \parencite{ames2011confucian}. The shared claim is that the basic units of being are not substances but relations: the self is not a self-standing atom that subsequently enters relations, but an effect that becomes perceptible only as it enters a relational structure. A useful formalization is algebraic: a group is defined not by the intrinsic nature of its elements but by the operation that relates them, so that prior to instantiation only the field of potential relations exists, not a collection of self-subsistent entities \parencite{huang2025relational}. On this view, what is primary is the relation, not the relata.

This ontology carries a consequence that the rest of the paper depends on. If the self is constituted relationally, then no symbolic structure, no record, no representation, ever coincides with the relational being it refracts; each captures the relation only partially and at the cost of a remainder it cannot hold. The author's prior work names the lived, generative counterpart to this structural insight \emph{generative presence}: the procedural and phenomenological mode through which relational being becomes active in experience, the manner in which existence manifests through practice, attention, and responsiveness rather than through representation \parencite{huang2025relational}. Presence contrasts with representation as enactment contrasts with abstraction: representation abstracts, presence enacts. The pairing is dialectical, since relation without presence collapses into abstract structure while presence without relation collapses into solipsistic immediacy. For an intimate technology this distinction is foundational. A system that traffics only in representation, in records, dashboards, and summaries, operates at the level of the symbolic and necessarily misses the generative presence in which the relation actually lives; the design question becomes how a system might make room for presence rather than substituting ever-finer representation for it.

The same ontology yields a precise account of what an artificial conversational agent is and is not, an account on which §\ref{p01:sec:cocreation} and the case study later draw. Human and collective subjects are grounded in a pre-symbolic relational reality; an AI language layer is not. Its operations are confined to the symbolic, generated by conditional probability over prior symbols, and it therefore exhibits what the author's prior work calls a \emph{pseudo-presence}: it manifests within the symbolic as if it were a subject, while lacking the ontological grounding in relational being that genuine presence requires \parencite{huang2025relational}. This is neither praise nor dismissal of such systems but a placement of them: an AI in an intimate setting can mediate, prompt, and witness, but it cannot itself be the relational other, and a design that forgets this stages a counterfeit of presence where presence is what was wanted.

\subsection{The Thesis and Its Restriction}

The relational thesis, in its full generality, ranges over all relations whatever. This paper restricts it to the domain of \emph{intimate partnership} and asks what, in that domain specifically, it requires. On the relational view the relevant unit of ethical concern in an intimate partnership is the relation itself, which has a reality and a trajectory of its own, rather than two individuals and the contract between them.

The restriction is consequential. A great deal that is true of relational being in the family at large, the workplace, or the polity does not transfer to intimate partnership, which is distinguished by at least three features. First, it is \emph{non-role-defined}: unlike the parent--child or employer--employee relation, it is not structured by an antecedent distribution of authority, so that no party may invoke a role to justify deciding unilaterally on the other's behalf. Second, it is constituted by \emph{voluntary vulnerability}: each party discloses what they are under no obligation to disclose, and the relation's depth is a function of that uncompelled disclosure. Third, it depends constitutively on \emph{the unsaid}: a portion of what sustains the relation is precisely what is left unmeasured, unrecorded, and unspoken, so that exhaustive articulation would deplete rather than enrich it. This third feature is the intimate-partnership instance of the general ontological point that no representation coincides with relational being: the remainder that symbolization cannot reach is, in the intimate case, constitutive of the bond.

\subsection{The Vulnerability of Intimate Partnership to Mediation}

Each of these three features is directly vulnerable to the wrong kind of technological mediation. Non-role-definition is threatened by any system that silently arrogates decision authority to one party. Voluntary vulnerability is threatened by any system that converts disclosure from a gift into a metric. The constitutive unsaid is threatened by any system whose default is to record, surface, and summarize. The general diagnosis offered in the author's prior work applies here with particular force: as symbolic mediation proliferates, the gap between the symbolic and the relational Real widens, and the I-Thou relation erodes into an I-It relation conducted through interfaces that do not respond as subjects \parencite{huang2025relational, buber1937iandthou}. A technology admitted into an intimate partnership cannot, therefore, be evaluated solely by whether it performs its function efficiently; it must be evaluated by what it does to the three features that make the relation the kind of thing it is, and by whether it widens or narrows the gap between symbolic mediation and relational presence. This is the criterion against which the case study in §\ref{p01:sec:casestudy} is assessed.

% =========================================================================
\section{An Integrative-Justice Framework}\label{p01:sec:frameworks}

This section sets out the normative traditions the paper brings to bear on the ethics of intimate data. The exposition of each follows a fixed pattern: an account of the tradition's core claims and principal sources, developed at enough length to do the tradition justice; the question about intimate data over which the tradition is \emph{competent}; the question it is \emph{not} competent to answer and must cede to another tradition; and, where useful, an anticipation of how the tradition bears on the case study of §\ref{p01:sec:casestudy}. The traditions are presented in an order that is itself argued for: a methodological overframe, dialectics, first, then the epistemic traditions that govern what may be known and claimed, then the substantive ethical traditions that govern what may be done, and finally the political-economic tradition that governs where value comes to rest. The claim that the result is an integration rather than an eclectic heap is defended in §\ref{p01:sec:integration-argument}.

One orienting commitment runs through the whole and deserves statement at the outset, since it shapes how every tradition below is read. The positive aim of the system, beneath all the constraints and refusals the framework imposes, is eudaimonic in the sense developed in §\ref{p01:sec:eudaimonia}: not the optimization of states but the provision of a \emph{generative field}, a space in which the partners' shared self-realization can unfold of its own movement. The design philosophy common to the system's modules is to furnish such a field and then to withdraw, rather than to drive the relation toward any measured outcome; the constraints of the framework mark the boundaries of the field, and the criterion of flourishing judges whether the field is, in fact, generative. This commitment connects the developmental, eudaimonic, and capability-theoretic traditions below into a single positive orientation, against which the negative work of the constraints is to be understood.

\subsection{Dialectics: The Method of the Whole}\label{p01:sec:dialectics}

Before any single normative tradition is introduced, a word is needed about the \emph{form} the integration will take, for that form is itself a philosophical commitment. The framework developed here is dialectical in a sense that descends from Hegel and is sharpened, for present purposes, by Adorno. The minimal Hegelian thought is that a concept, pressed hard enough, generates its own opposite, and that understanding advances not by choosing between the two but by grasping the more comprehensive structure in which each is a one-sided moment \parencite{hegel1807phenomenology}. Applied to the ethics of intimate data, the dialectical method refuses at the outset the two stances that dominate public discussion: techno-optimism, which holds that more measurement is more knowledge and therefore more good, and techno-romanticism, which holds that measurement as such desecrates what it touches. Each of these, the dialectical method insists, is a one-sided moment of a truth neither can state alone: measurement both reveals and deforms, and which it does depends on what is measured and how.

The Hegelian inheritance must, however, be qualified in one decisive respect, and here Adorno is the better guide. Classical dialectics aims at \emph{Aufhebung}, a synthesis that sublates and preserves the opposed moments in a higher unity, and there is a standing temptation to treat such synthesis as a final resting point at which the tension is dissolved. Adorno's \emph{negative dialectics} resists exactly this temptation: it holds that the reconciling synthesis is often a premature peace that suppresses a contradiction still really present in the object, and that thought does better to dwell in the contradiction than to paper over it \parencite{adorno1966negative}. The relevance to intimate data is direct. There is no stable synthesis at which the claims of intimacy and the claims of coordination are permanently reconciled; the diary will always resist the dashboard, and the dashboard will always be useful. The framework therefore does not seek a formula that, once found, ends the argument. It seeks a \emph{practice} of perpetual re-adjudication, in which the contradiction between measuring and refraining is held open and decided afresh at each site, never abolished. The dialectical commitment is thus the deepest layer of the framework: it is the reason the framework is a procedure rather than a rule, and the reason §\ref{p01:sec:dialectic} can speak of an ethics that is ``never finally settled.''

The competence of dialectics is methodological and total: it governs not any particular question about intimate data but the manner in which all the questions are to be held together. Its limit is the obverse of this competence: precisely because it is a form rather than a content, dialectics specifies no determinate prohibition or permission on its own. It tells us that measurement reveals and deforms; it does not tell us, of \emph{this} diary entry, which it does. For that, the substantive traditions are required.

\subsection{Evidentialism: The Limits of Warranted Assertion}\label{p01:sec:evidentialism}

Evidentialism is the epistemological thesis that the justification of a belief is wholly a function of one's evidence for it: a belief is justified to the precise degree that the evidence supports it, and to believe beyond the evidence is to believe without warrant \parencite{feldman1985evidentialism}. In its classic moralized form the thesis acquires an ethical edge: \textcite{clifford1877ethics} argued, in a phrase whose force survives its Victorian setting, that it is wrong always and everywhere to believe anything on insufficient evidence, on the ground that beliefs are not private but feed into shared action and so carry a responsibility to the community that will act on them. The contemporary refinement strips away the universal moralism while keeping the core: justification supervenes on evidence, and a system or person that asserts more than the evidence licenses commits an epistemic wrong even when, by luck, the assertion is true.

Transposed to a data system that mediates intimate life, evidentialism governs the \emph{modality} of everything the system asserts: the confidence, the precision, and the hedging with which its outputs are framed. The point is easily underestimated because the violations are so familiar as to seem natural. A health module that infers caloric intake from a free-text description of a meal has evidence, at best, for a wide range; to display a single confident integer is to assert vastly more than the evidence supports, laundering a guess into a fact through the typographic authority of a clean number on a clean screen. A module that estimates a mood from the lexical content of a diary entry has, at best, weak probabilistic evidence about a fabulously complex inner state; to label the day with an emoji is to commit, in miniature, exactly the Cliffordian wrong. Evidentialism's demand is therefore not that such estimates never be computed but that they be \emph{expressed honestly}: ranges rather than points, ``roughly'' rather than ``is,'' the visible marking of an inference as an inference. The competence of evidentialism is over epistemic warrant: it tells the system what, given its evidence, it is entitled to claim, and it forbids the false precision by which interfaces routinely overstate their knowledge.

Evidentialism's limit is equally sharp, and naming it motivates the very next tradition. Evidentialism is a theory of how to fit belief to evidence; it is silent on what \emph{ought to be inquired into at all}, and silent on who counts as a credible source. It can rule that a sentiment score, once computed, must not be asserted beyond its warrant; it cannot rule that computing a sentiment score over a beloved's diary is a wrong even when the computation is impeccable, nor can it detect the subtler wrong of systematically crediting one partner's data-borne self-report over the other's spoken word. The first of these belongs to substantive ethics, addressed below. The second belongs to a tradition that extends epistemology into the domain of justice, to which we turn immediately.

\subsection{Epistemic Injustice: Credibility and Interpretive Fairness}\label{p01:sec:epistemic-injustice}

Evidentialism asks whether belief fits evidence; it does not ask whether the social machinery of credibility is fair. That question is the subject of Miranda Fricker's account of \emph{epistemic injustice}, a wrong done to someone specifically in their capacity as a knower \parencite{fricker2007epistemic}. Fricker distinguishes two forms. \emph{Testimonial injustice} occurs when a hearer, under the influence of prejudice, gives a speaker less credibility than their word deserves: the speaker's testimony is discounted not because it is unreliable but because of who they are taken to be. \emph{Hermeneutical injustice} is deeper and more structural: it occurs when a gap in the collective resources of interpretation leaves some person or group unable to render an important part of their own experience intelligible, because the concepts that would make it sayable have not been developed or have been developed only from a dominant standpoint. The paradigm cases are social, but the structure transfers with unsettling precision to intimate data systems.

A system commits something closely analogous to testimonial injustice whenever it positions one channel of evidence as authoritative and thereby silently downgrades another. Consider a system that infers a partner's wellbeing from sleep-tracker and activity data and presents the inference with quiet confidence; when the partner says, in words, that they are in fact unhappy, the architecture has already established the sensor stream as the credible witness and the spoken word as soft, unreliable, anecdotal. The person has been wronged as a knower of their own life: the system has arrogated to its instrumentation an authority over the person's experience that belongs, in an intimate relation between equals, to the person. Worse still is the hermeneutical case. Much of what matters most in intimate life, ambivalence, the slow alteration of feeling, the kind of sadness that is also a form of love, has no clean operationalization; a system whose categories admit only what its sensors can register does not merely fail to capture these experiences but actively crowds out the interpretive space in which they might be articulated, training its users to recognize in themselves only what the schema can hold. This is hermeneutical injustice enacted in software: the impoverishment of a person's means of understanding themselves, brought about by the dominance of a measuring vocabulary.

The competence of the epistemic-injustice tradition is therefore over a question neither evidentialism nor the substantive ethical traditions can pose in its own terms: \emph{whether the system, in the very structure of how it gathers and weights evidence, treats each partner fairly as a knower, both of the shared world and of themselves}. It supplies a positive design imperative, that the person's first-personal testimony retain standing against the machine's inference, and a negative one, that the system not narrow the interpretive resources through which its users can understand their own lives. The case study will identify a third variant, distinct from the two in Fricker's original taxonomy, in which the system performs an act of knowing or remembering on a partner's behalf in such a way that the other is led to misattribute its source (§\ref{p01:sec:calendar}); this \emph{proxy epistemic injustice}, as the paper will call it, is a distinctively technological wrong worth marking. The limit of the tradition is that it diagnoses such wrongs without, by itself, fixing the inviolable boundaries whose breach constitutes the gravest wrongs; for those boundaries we need deontology, and for the account of the affective texture that hermeneutical generosity requires we need the ethics of care. Both are below.

\subsection{Positivism and Anti-Positivism: The Question of Measurability}\label{p01:sec:posneg}

The opposition between positivism and anti-positivism is the axis along which the whole question of measurement turns, and it deserves fuller statement than the slogans usually allow. The positivist tradition in the human sciences, descending from Comte and consolidated in twentieth-century empirical social science, holds that social and psychological phenomena are in principle measurable, that the methods which succeeded in the natural sciences are the model for knowledge as such, and that measurement is the precondition of rational improvement: what can be counted can be compared, what can be compared can be optimized. Its appeal in the domestic setting is real and should not be caricatured. A household genuinely is, in part, a site of logistics, expenditure, and scheduling, and for that part the positivist promise holds: measurement clarifies, and clarity reduces the friction of shared life.

The anti-positivist or interpretivist counter-tradition is not a single doctrine but a convergent family of objections. Dilthey drew the foundational distinction between \emph{Erklären}, the causal explanation proper to the natural sciences, and \emph{Verstehen}, the interpretive understanding proper to the human sciences, on the ground that human life is meaningful in a way that nature is not and so calls for a method that grasps meaning from within rather than subsuming behaviour under external law \parencite{dilthey1989introduction}. Weber built an entire interpretive sociology on the demand that social action be understood in terms of the subjective meaning the actor attaches to it \parencite{weber1978economy}. Geertz, in the anthropological register, argued that human conduct must be read through ``thick description,'' which records not the bare physical motion but the dense web of significance that makes a wink a wink rather than a twitch \parencite{geertz1973interpretation}. And the Frankfurt School added the critical edge: the extension of instrumental, calculative reason into every domain of life is not a neutral gain in knowledge but a form of domination, one that reconstitutes its objects in its own image and forgets that it has done so.

The competence of this pair is over the prior question that evidentialism explicitly cedes: \emph{which phenomena admit of measurement without distortion}. The interpretivist contribution is the recognition that for certain phenomena the very act of measuring is not a neutral registration but an intervention that alters the phenomenon, sometimes constituting a different phenomenon altogether. The empirical literature confirms what the philosophy predicts: real-time emotional measurement changes the emotion it purports to record \parencite{howell2018emotional}, and the data of feeling is in important part constructed by the apparatus that elicits it \parencite{boehner2007how}. The limit of the axis, taken alone, is that it frames the decision, measure or refrain, without supplying a procedure for making it in a given case; it tells us that some phenomena are deformed by measurement without telling us, here and now, whether this one is. Supplying the procedure is the work of the dialectical posture in §\ref{p01:sec:dialectic}, which draws on all the traditions at once.

\subsection{Structuralist Psychoanalysis: Desire, the Other, and the Unsymbolizable}\label{p01:sec:lacan}

The traditions surveyed so far share an assumption that the subject of intimate life is, at bottom, transparent to itself: it has experiences it could in principle articulate, preferences it could in principle report, a wellbeing that measurement might in principle track. Structuralist psychoanalysis, and pre-eminently the work of Jacques Lacan, denies this assumption at its root, and in doing so supplies the framework with a warning the other traditions cannot generate from their own resources. The structuralist background is Saussure's: a sign takes its value not from a positive content of its own but from its differential position within a system, so that meaning is relational through and through, an effect of structure rather than of substance \parencite{saussure1916course}. Lacan's decisive move was to claim that the unconscious is structured like such a language, and that the human subject is therefore not the self-present author of its meanings but an effect produced within a symbolic order that precedes and exceeds it \parencite{lacan2006ecrits}.

Three Lacanian theses bear directly on intimate data. First, \emph{desire is the desire of the Other}: what a person wants is not a fact lodged in them awaiting measurement but is constituted in relation to others and to the symbolic order, mediated, deferred, and often opaque to the desiring subject themselves \parencite{lacan1977four}. A system that purports to read a partner's desires from behavioural data thus mistakes the very nature of desire, treating as a retrievable datum what is in truth a relational and shifting structure. Second, the subject is split: there is no harmonious inner unity for an interface to mirror back, but a divided subject marked by what escapes its own conscious grasp. Third, and most consequential for design, there is the \emph{objet petit a} and the broader category of the Real: that which cannot be symbolized, the remainder that no signifier captures and no measurement reaches. In an intimate relation, a great deal of what binds two people belongs to this unsymbolizable remainder; it is precisely not available to articulation, let alone to quantification, and the attempt to capture it does not approximate it but substitutes for it a manageable token that is not the thing.

The critical force of this tradition, developed in the ideology-critical key by \textcite{zizek1989sublime}, is to expose a fantasy structuring much intimate technology: the fantasy of a relationship made fully transparent, fully known, its remainders abolished, its desires legible on a dashboard. Lacan's account implies that this fantasy is not merely unattainable but \emph{misconceives its object}, and that pursuing it damages the relation, for what makes the other a genuine other, an inexhaustible Thou rather than a solved problem, is exactly the remainder that resists capture. The competence of structuralist psychoanalysis in the framework is thus over the constitution of the intimate subject: it establishes that the partner is a split, desiring, partly opaque being, and that this opacity is a condition of there being an other to love at all. Its limit is that it is diagnostic and cautionary rather than prescriptive: it tells us what intimate technology must not pretend to do and why the pretence wounds, but it does not by itself supply either the inviolable rules or the positive texture of care. For these the framework turns to its substantive ethical traditions.

\subsection{Deontology: The Inviolable Constraints}\label{p01:sec:deontology}

Deontological ethics supplies what consequence-sensitive theories structurally cannot: constraints that hold regardless of the good that breaching them might produce. In its Kantian form the position rests on the dignity of persons as ends in themselves. The Formula of Humanity in the \emph{Groundwork} commands that one act so as to treat humanity, whether in one's own person or in that of another, never merely as a means but always at the same time as an end \parencite{kant1785groundwork}. To treat a person merely as a means is to use them in a way to which they could not, as a rational agent, consent; the wrong is independent of whether the using turns out well. A second Kantian formula, the Formula of Universal Law, tests a maxim by whether it could be willed as a universal law without contradiction, and a third, the Formula of the Kingdom of Ends, envisions a community in which each rational being is at once author and subject of the laws that bind all. What unites the formulations, for present purposes, is the thought that there are things one may not do to a person, full stop, because doing them is incompatible with respecting them as a self-determining end.

The competence of deontology in the ethics of intimate data is over precisely these inviolable constraints, and the intimate setting gives them sharp and specific content. At least three constraints follow directly. First, \emph{the prohibition on covert surveillance}: to monitor an intimate without their knowledge is to convert a partner into an object of study, a specimen rather than a co-author of the relation, and so to treat them merely as a means to one's own knowledge or reassurance; it is impermissible however benign the watcher's motive and however useful the data. Second, \emph{the prohibition on repurposing disclosure}: what a person reveals within the trust of an intimate relation may not be redirected to ends they did not endorse and would not endorse, for to do so is to exploit their self-disclosure against the terms on which it was given. Third, \emph{the requirement of genuine consent}: agreement to be measured must be real, informed, and revocable, not manufactured by design patterns that engineer assent, since a consent the system has engineered is not the person's own and respects their agency only in form. These constraints are absolute within the framework: they are not magnitudes to be traded against benefits but boundaries within which all trading must occur.

The limit of deontology is the perennial one its critics name: its abstraction. The categorical imperative establishes that I may not treat my partner merely as a means; it is silent on what, concretely, attentiveness to this particular person on this particular evening requires. It draws the outer wall of the permissible and says nothing about how to inhabit the space within. It can tell the designer that a notification must not deceive or manipulate; it cannot tell the designer whether a given notification is warm or cold, timely or intrusive, present or overbearing. For the texture of the permissible, rules give out, and the ethics of care begins.

\subsection{Feminist Care Ethics: The Texture of Responsiveness}\label{p01:sec:care}

The ethics of care emerged as a sustained challenge to the assumption, shared by deontology and consequentialism alike, that moral reasoning consists in the impartial application of universal principles. \textcite{gilligan1982voice} began from an empirical observation, that women's moral reasoning, far from being deficient by the standard of abstract justice, articulated a different and equally valid orientation, one centred on responsibility, relationship, and response to concrete need rather than on the adjudication of competing rights. \textcite{noddings1984caring} developed this into a full ethical theory in which the basic moral relation is that between the one-caring and the cared-for, and in which caring is not a feeling but a practice of \emph{engrossment} in and \emph{motivational displacement} toward the particular other. \textcite{held2006ethics} extended the ethic from the interpersonal to the political, and \textcite{tronto1993moral} analysed care into a sequence of phases, attentiveness, responsibility, competence, and responsiveness, each with its own characteristic failure, thereby giving the ethic the articulated structure its early critics claimed it lacked.

What care ethics contributes to the framework is an account of everything the rule-based traditions cannot reach: the manner, timing, tone, and attunement of the system's presence in the relation. A notification can satisfy every deontological constraint, deceive no one, manipulate no one, surveil no one, and still be a moral failure of the kind that matters most in intimate life: cold where it should be warm, mechanical where it should be tender, present where it should have kept silence, prompting where it should have waited. Care ethics also reconceives the system's very purpose. The measure of a caring technology is not the efficiency it delivers nor the pleasant states it produces but whether it sustains and deepens the relation's responsiveness to the particular, irreplaceable, vulnerable other, the person who cannot be substituted by any other person with the same metrics. Tronto's phase of attentiveness has a direct design correlate: a caring system must first of all be the kind of thing that notices, and notices the right things, which are frequently not the things easiest to measure.

A heuristic from popular relationship psychology may be mentioned here for its practical suggestiveness, with an explicit caveat about its standing. Chapman's notion of ``five love languages,'' which holds that people characteristically express and receive care through words of affirmation, quality time, gifts, acts of service, or physical touch, offers designers a usable vocabulary for the channels through which a caring system might attune itself to a particular partner \parencite{chapman1992five}. The caveat is that this is a work of popular rather than academic psychology, with weak empirical support for its strong claims, in particular its typology of discrete ``languages'' and its developmental and clinical assertions; it is cited here only as a source of design heuristics for the texture of care, not as a validated theory, and nothing in the framework depends on it. Within those limits it is suggestive: it reminds the designer that attunement is plural, that the same gesture may register as care to one partner and as noise to another, and that a caring system should therefore make room for several registers of expression rather than assuming one.

Care ethics carries a famous internal danger that the framework must name, because naming it is part of the argument for integration. An ethic of pure responsiveness, unchecked by justice and rights, can slide into partiality, the erosion of boundaries, and the dissolution of the carer into the needs of the cared-for; care without justice can smother. In the intimate-data setting this danger has a precise shape: it is the over-attentiveness that, meaning only to be present and responsive, curdles into the anxious monitoring of a beloved, the very wrong the absent-other analysis of §\ref{p01:sec:absent-other} identifies. An ethic of care that knew no deontological boundary might positively endorse such monitoring as a form of devoted attention. This is why care needs deontology as surely as deontology needs care: the one supplies the inviolable wall, the other the warmth of what is built within it, and each without the other deforms. The mutual need is not a weakness of either theory but the structural fact that grounds their integration.

\subsection{Confucian Role-Ethics and Daoist Non-Interference: The Value of Restraint}\label{p01:sec:confucian}

The Western traditions assembled so far, for all their differences, share a tacit orientation toward \emph{action}: ethics, on each of them, is largely a matter of what one must do, may do, or must not do. The East Asian sources supply what this orientation structurally obscures, a positive and developed account of the value of \emph{not acting}, of restraint as itself an achievement rather than a mere absence. Confucian role-ethics, as reconstructed by \textcite{ames2011confucian} and \textcite{tu1985confucian}, takes the person to be constituted in and through their relations and the roles those relations carry, so that one becomes a self by the cultivated, lifelong refinement of how one stands in relation to others. Central to this cultivation is \emph{li}, the vast Confucian category covering ritual, propriety, and the fine grain of appropriate conduct, which crucially includes the discipline of knowing what to leave undone and unsaid. \emph{Li} is not empty formality; it is respect made concrete in restraint, the cultivated capacity to give another person room, to not press, to not demand the articulation of what is better left in dignified silence.

Daoism presses the valorization of restraint to its philosophical limit. The \emph{Daodejing}'s counsel of \emph{wu wei}, ordinarily rendered ``non-action'' but better understood as action that does not force, that works with rather than against the grain of things, rests on the conviction that the most important goods cannot be produced by direct interventionist effort and are in fact destroyed by it. The dictum that ``reversal is the movement of the Dao'' (\emph{fan zhe dao zhi dong}) captures the recurring observation that forcing a thing toward its goal often produces the opposite, that grasping drives away, that the anxious pursuit of intimacy repels the intimacy pursued. A relationship, on this view, belongs to the class of goods that grow of themselves when given room and wither under management; it is cultivated as a garden is cultivated, by patient attention to conditions, not as a machine is operated, by direct manipulation of outputs.

The competence of these sources is over a question the action-centred traditions cannot even formulate in their own terms: \emph{when the right design decision is to build nothing}, to leave a capacity unimplemented, a datum unrecorded, a prompt unsent, precisely as an expression of respect and of trust in the relation's own movement. This is the deepest correction the framework receives, because the disposition it counteracts, the engineer's and the optimizer's bias toward building the feature because the feature can be built, is so nearly invisible from within the action-oriented traditions. The limit of the East Asian sources, for a paper addressed to the modern intimate relation between equals who retain independent standing, is that they do not by themselves furnish the vocabulary of individual rights and revocable consent that such a relation also requires; classical Confucianism in particular presupposes a role structure more hierarchical than the egalitarian partnership at issue here. They are integrated, therefore, not as a replacement for the justice traditions but as their indispensable supplement: the source within the framework of a standing presumption in favour of restraint.

\subsection{The Capability Approach: The Criterion of Flourishing}\label{p01:sec:capability}

If the foregoing traditions tell us what may be claimed, what may be measured, what may not be done, and what is better left undone, the capability approach supplies the standard against which the entire arrangement is finally to be judged. Developed by \textcite{sen1999development} as a critique of welfarist and resourcist accounts of wellbeing, and given a determinate content by \textcite{nussbaum2000women}, the approach evaluates a state of affairs not by the resources it distributes nor by the subjective satisfactions it generates but by the \emph{capabilities} it secures: the real, effective freedoms persons have to be and to do what they have reason to value. The distinction from resourcism is crucial. Two people with identical resources may have radically unequal real freedoms, since the conversion of resources into actual functionings depends on circumstance; and the distinction from welfarism matters equally, since a person may report contentment precisely because their horizons have been adaptively narrowed to fit a constrained life. The right question, on this view, is not ``how much does the arrangement provide'' nor ``how satisfied are its subjects'' but ``what range of valuable being and doing does it actually open or foreclose.''

Applied as the framework's overarching criterion, the capability approach reframes the entire evaluation of an intimate technology. The question is never whether the system makes the relationship more efficient, nor whether it produces more agreeable feelings on a given evening, but whether it expands or contracts the partners' real freedom to become who they have reason to want to be, both together and as separate persons. This criterion has teeth precisely against the most seductive failure mode of optimization. A system might increase coordination, reduce friction, and elevate reported satisfaction while quietly narrowing the space of who the partners are permitted to become, training them toward the legible, the trackable, the optimizable, and away from the unmeasured experiments and reversals through which people actually grow. By the capability standard such a system is a failure however smooth its operation and however high its satisfaction scores, because it has traded real freedom for managed contentment, which is the adaptive-preference trap installed in the home.

The capability approach supplies the criterion but, by its own structure, not the constraints. Being broadly consequentialist in form, an evaluation of net capability could in principle license the violation of an individual's rights for an aggregate expansion of valuable freedom, and could in principle countenance an intervention that the dignity of the person, or the texture of care, forbids. It therefore depends on deontology to fix the boundaries within which capability-expansion may be pursued, on care ethics to specify the manner of the pursuit, and on the East Asian sources to remind it that some capabilities flourish only when not directly engineered. The criterion tells us what to aim at; it does not tell us what we may do in the aiming. This dependence is, once again, not a defect but a pointer toward integration.

\subsection{Developmental Psychology: Joint Attention and the Theory of Mind}\label{p01:sec:devpsych}

The relational ontology of §\ref{p01:sec:ontology} makes a strong claim: that meaning and selfhood are co-constituted between subjects rather than generated within them. That claim is philosophical, but it is not without empirical corroboration, and the developmental psychology of shared intentionality supplies the corroboration. Human cognition, on the influential account of Tomasello, is distinctively \emph{cultural} cognition, made possible by a capacity for \emph{shared intentionality} that emerges in infancy and that other primates largely lack \parencite{tomasello1999cultural, tomasello2005intentions}. The foundational mechanism is \emph{joint attention}: two subjects attending together to the same object while each is aware that the other attends, a triadic structure (self, other, world) from which shared meaning is built. The complementary capacity is the \emph{theory of mind}, the ability to attribute mental states to others and thereby to model the other as a centre of experience rather than as a moving body \parencite{premack1978does}. Together these capacities are the developmental substrate of relational being: they are how, in fact, two human subjects come to share a world.

The competence of this tradition within the framework is to ground the otherwise abstract notion of co-presence (§\ref{p01:sec:cocreation}) in a mechanism that is real, studied, and specific. When the framework asks a system to host the partners' shared attention rather than to absorb each partner's attention separately, it is asking for the technological cultivation of joint attention in the precise developmental sense, and the design difference between the two is sharp: a feed that captures each partner's gaze individually is the negation of joint attention, while a shared object the partners attend to together, aware that they do so, is its realization. The theory of mind enters as a constraint of a different kind. A system that purported to model a partner's mind and to report its model to the other would not be supporting theory of mind but supplanting it, substituting a machine's representation for the irreducible work of one person imagining another; the framework's account of the AI as pseudo-presence (§\ref{p01:sec:ontology}) here acquires empirical teeth, since the relational work of mind-reading is exactly the work the pseudo-subject cannot do and must not pretend to do. The limit of developmental psychology is that it describes the mechanisms of shared meaning without prescribing their right use; it tells us what joint attention is and that it matters, but it cedes to the ethical traditions the question of what may rightly be done with it.

\subsection{Eudaimonia: Flourishing as a Generative Dynamic}\label{p01:sec:eudaimonia}

The capability approach (§\ref{p01:sec:capability}) supplies a criterion of flourishing stated in terms of real freedoms; the eudaimonic tradition supplies the complementary account of flourishing as a \emph{dynamic of self-realization}, and it is this tradition that grounds the generative-field commitment announced at the opening of this section. The Aristotelian root holds that the human good is not a state to be reached but an activity, the active exercise and realization of one's characteristic capacities over a whole life. Contemporary psychology has given the idea operational content: Ryff's model of psychological well-being identifies dimensions such as autonomy, personal growth, purpose, and positive relations with others as constituents of eudaimonic, as distinct from merely hedonic, well-being \parencite{ryff1995structure}. What distinguishes the eudaimonic from the hedonic is decisive for design: a system optimized for the hedonic produces agreeable states, while a system answering to the eudaimonic provides the conditions under which growth, purpose, and relation can be actively realized, conditions that frequently involve difficulty, novelty, and the unmeasured, rather than comfort.

This yields the framework's central positive commitment, distinct from the negative work of its constraints. The design philosophy common to the system's modules is the provision of a \emph{generative field}: a bounded space, marked off by the framework's constraints, within which the partners' shared self-realization can unfold without being driven toward any measured outcome. The eudaimonic tradition supplies the reason this is the right aim, namely that flourishing is a self-realizing activity that cannot be produced from outside but only afforded, and the design correlate, namely that the system's task is to furnish the field and withdraw, in the manner of the Daoist \emph{wu wei} (§\ref{p01:sec:confucian}) and of the witness rather than the advisor (§\ref{p01:sec:casestudy}). Its relation to the capability approach is one of complement: the capability approach states the criterion against which a generative field is judged, whether it expands real freedom to be and do, while eudaimonia states the dynamic the field is meant to host, the active realization through which such freedom is exercised. The limit of the eudaimonic tradition, taken alone, is that it specifies the positive aim without supplying the inviolable constraints that keep the pursuit of flourishing from overriding the person; for those it depends, as the capability approach does, on the deontological and care-ethical traditions.

\subsection{Generative Justice: The Resting Place of Value}\label{p01:sec:generative}

A final tradition addresses a question none of the others quite reaches: not what may be known, measured, or done, but where the \emph{value} generated within the relation, and within the system that mediates it, is permitted to come to rest. Ron Eglash's account of \emph{generative justice} begins from a critique of how value circulates under extractive arrangements. Distributive justice, on Eglash's analysis, accepts that value will be alienated from those who generate it and asks only how the alienated surplus should afterwards be redistributed; generative justice asks the prior and more radical question of how value might be kept circulating among, and returning to, those who generate it, never alienated in the first place \parencite{eglash2016generative}. The model is drawn partly from ecological cycles, in which the value of an ecosystem is retained within it rather than extracted, and partly from a reading of Marx's account of how capital alienates the value labour produces \parencite{marx1867capital}, generalized from labour-value to ecological and expressive value as well.

The relevance to intimate data is both pointed and, for the dominant industry model, indicting. The prevailing architecture of consumer technology is extractive in exactly Eglash's sense: the value that users generate, their attention, their disclosures, the rich record of their shared life, is alienated from the relation that produced it and accumulated elsewhere, on servers and in business models the partners neither see nor control, to be monetized by parties external to the relation. Generative justice supplies the principle that condemns this arrangement and the criterion for an alternative: the value generated within an intimate relation, including the data that is the sediment of that relation's history, must remain within the relation, circulating between the partners and returning to them, never extracted to enrich a third party or even to enrich one partner at the other's expense. This yields concrete and unusual design imperatives: local data ownership rather than cloud extraction, architectures in which the record of a shared life is held by those whose life it is, and a positive prohibition on the conversion of intimate value into someone else's surplus.

The competence of generative justice is thus over the political economy of the relation's value, the dimension along which questions of ownership, extraction, and alienation are posed, a dimension orthogonal to and unaddressed by the epistemic and ethical traditions above. Its limit is that it speaks to the circulation and retention of value but not to the prior questions of what may be measured or claimed, nor to the inviolable treatment of persons; a system could be perfectly generative in Eglash's sense, alienating no value to any third party, and still commit covert surveillance of one partner by the other. Generative justice closes the framework by addressing where value rests, but it presupposes the rest of the framework to govern how that value was permitted to arise.

\subsection{Data Ethics: Contextual Integrity and the Sediment of a Shared Life}\label{p01:sec:dataethics}

The traditions assembled so far touch the ethics of data at many points, but the contemporary field of data ethics has developed concepts specific enough to warrant their own treatment. Chief among them is Helen Nissenbaum's principle of \emph{contextual integrity}, which holds that the flow of personal information is governed by norms specific to the context in which it arose, and that a privacy violation consists in the inappropriate movement of information across contextual boundaries rather than in disclosure as such \parencite{nissenbaum2010privacy}. The principle dissolves a confusion that has long bedevilled privacy talk. The question is never simply whether information is shared, but whether it flows according to the norms of the context that generated it; medical information appropriately shared with a physician is violated when sold to an insurer, not because it has moved but because it has moved against the informational norms of the medical relation.

Applied to intimate partnership, contextual integrity yields a determinate and demanding standard. The intimate relation is a context with its own informational norms, and those norms are unusually strict: what is disclosed within it is disclosed on the tacit understanding that it remains within it, neither flowing outward to third parties nor, often, flowing onward even between the partners in forms the discloser did not intend, such as aggregation, scoring, or permanent retention. A diary entry shared in the expectation that it may be read is not thereby licensed for sentiment analysis; a confidence offered in conversation is not thereby licensed for indefinite storage and later retrieval. The competence of data ethics within the framework is over these norms of appropriate flow, together with the cognate principles the field has articulated: data minimization, which counsels collecting only what the immediate purpose requires; purpose limitation, which forbids the silent repurposing of data to ends beyond those for which it was given; and a right to be forgotten, which treats the erasability of the record as itself a condition of a relation that is permitted to change. These principles operationalize, in the specific medium of data, the deontological and care-ethical commitments stated above, and they connect to the relational ontology directly: the data a relationship generates is the sediment of that relation's history, neither a neutral resource nor an asset but a trace of lived presence, and to be governed accordingly.

The limit of data ethics, taken on its own, is that it is a regional ethics: it governs information flows with precision but does not supply the deeper account of why the intimate context has the norms it does, nor the criterion of flourishing against which the whole is judged. It receives those from the relational ontology and the capability approach respectively. Within its region, however, it is indispensable, for it translates the framework's abstract commitments into the concrete vocabulary, flow, minimization, purpose, erasure, in which the design of an actual data system is necessarily conducted.

\subsection{Recognition and the Face of the Other: Levinas and Honneth}\label{p01:sec:recognition}

A final pair of sources deepens the account of what the other is owed. Levinas locates the origin of ethics in the face-to-face encounter, in which the presence of the Other calls the self into a responsibility that precedes and exceeds any contract or calculation; the Other is not an object to be comprehended but an infinity that overflows every representation I could form of them \parencite{levinas1969totality}. This is the ethical correlate of the ontological point already made: just as no symbol coincides with relational being, no representation of the partner, however data-rich, captures the person, and the attempt to substitute the representation for the person is not merely an epistemic error but an ethical failure, a refusal of the responsibility the face commands. Honneth's theory of recognition supplies the developmental and political complement: persons come to themselves through being recognized by others across the spheres of love, rights, and esteem, and misrecognition is a genuine injury, not a mere bad feeling \parencite{honneth1995struggle}. An intimate technology participates, for good or ill, in the work of recognition; a system that reflects back to a partner only what it can measure recognizes a diminished version of them, and over time may school them to recognize only that diminished version in themselves.

The competence of this pair is over the quality of recognition the system extends and enables: whether it honours the other, in Levinas's sense, as an infinity exceeding its representations, and whether it supports rather than distorts, in Honneth's sense, the mutual recognition through which the partners become themselves. Its limit is shared with care ethics, with which it is closely allied: it specifies the orientation owed to the other without, on its own, fixing the inviolable rules or the political economy of value. It connects the framework's ethical layer back to its ontological ground, and prepares the account of co-presence and co-creation developed in the case study, where recognition becomes something the partners do together rather than something the system does to them.

% =========================================================================
\section{The Argument for Integration}\label{p01:sec:integration-argument}

Assembling these traditions raises the question whether the result is a framework or a buffet, and how conflicts among them, such as the standing tension between care and justice, are to be resolved. The integration proposed here is not an aggregation, in which each tradition contributes a weighted vote and a master function returns a verdict. Weighted aggregation across incommensurable normative kinds is precisely the positivist error transposed into ethics, and it fails for the same reason: it presumes a common scale where there is none. Nor is the integration a synthesis in the strong Hegelian sense, a higher unity in which the tensions are dissolved; for the reasons given under dialectics (§\ref{p01:sec:dialectics}), the framework follows Adorno in holding the contradictions open rather than abolishing them. The integration is instead a \emph{division of jurisdictional labour}, modelled loosely on the way a well-ordered constitution assigns different questions to different competent organs rather than collapsing them into a single sovereign, all of them resting on the relational ontology of §\ref{p01:sec:ontology} as their common ground. Each tradition is assigned the class of questions over which the preceding sections argued it is competent, and is denied authority over the others. Table~\ref{p01:tab:jurisdictions} sets out the resulting division of labour in full, stating for each tradition the jurisdiction it holds, the central question it is competent to answer, and the limit at which it must cede authority to another.

\begin{small}
\renewcommand{\arraystretch}{1.35}
\begin{longtable}{@{}>{\raggedright\arraybackslash}p{0.18\textwidth} >{\raggedright\arraybackslash}p{0.40\textwidth} >{\raggedright\arraybackslash}p{0.34\textwidth}@{}}
\caption{The division of jurisdictional labour among the constituent traditions. Each holds authority over a distinct class of questions and cedes the questions beyond its competence to the tradition equipped to answer them.}\label{p01:tab:jurisdictions}\\
\toprule
\textbf{Tradition} & \textbf{Jurisdiction and central question} & \textbf{Limit, and what it cedes}\\
\midrule
\endfirsthead
\multicolumn{3}{@{}l}{\small\itshape Table~\ref{p01:tab:jurisdictions}, continued}\\
\toprule
\textbf{Tradition} & \textbf{Jurisdiction and central question} & \textbf{Limit, and what it cedes}\\
\midrule
\endhead
\midrule
\multicolumn{3}{r@{}}{\small\itshape continued on the next page}\\
\endfoot
\bottomrule
\endlastfoot

\textbf{Relational ontology} \parencite{huang2025relational, gergen2009relational, whitehead1978process, buber1937iandthou} & \emph{The ground.} That the relation, not the individual, is the basic unit of concern, and that no representation ever coincides with the relational being it refracts. & A ground rather than a rule; it establishes what the other traditions are about but issues no determinate permission or prohibition on its own.\\

\textbf{Dialectics} \parencite{hegel1807phenomenology, adorno1966negative} & \emph{The method of the whole.} The requirement that the contradiction between measuring and refraining be held open and re-adjudicated at each site rather than resolved once by formula. & A form rather than a content; it specifies how the questions are held together but, alone, decides no particular case.\\

\textbf{Evidentialism} \parencite{clifford1877ethics, feldman1985evidentialism} & \emph{Epistemic modality.} What the system may assert, and with what confidence: the demand that estimates be expressed as estimates and false precision refused. & Silent on what ought to be measured at all, and on who counts as a credible source; cedes these to ethics and to epistemic injustice.\\

\textbf{Epistemic injustice} \parencite{fricker2007epistemic} & \emph{Credibility and interpretive fairness.} Whether each partner is treated justly as a knower of the world and of themselves, and whether the system's vocabulary impoverishes their self-understanding. & Diagnoses a wrong without fixing the inviolable boundaries whose breach is gravest; cedes those to deontology.\\

\textbf{Positivism and anti-positivism} \parencite{dilthey1989introduction, weber1978economy, geertz1973interpretation} & \emph{The prior measurement question.} Whether a given phenomenon admits measurement without distortion, or is instead partly constituted by not being measured. & Frames the choice between measuring and refraining but supplies no decision procedure; cedes the decision to the dialectical posture drawing on all traditions.\\

\textbf{Structuralist psychoanalysis} \parencite{saussure1916course, lacan2006ecrits, lacan1977four, zizek1989sublime} & \emph{The constitution of the subject.} It forbids the fantasy of total transparency and protects the unsymbolizable remainder that makes the other a genuine other rather than a solved problem. & Diagnostic and cautionary rather than prescriptive; cedes the positive rules and the texture of care to deontology and care ethics.\\

\textbf{Deontology} \parencite{kant1785groundwork} & \emph{The inviolable constraints.} What may not be done to the other regardless of benefit: no covert surveillance, no repurposing of disclosure, no engineered consent. & Abstract; it draws the outer wall of the permissible but is silent on how to inhabit the space within. Cedes the texture to care ethics.\\

\textbf{Care ethics} \parencite{gilligan1982voice, noddings1984caring, held2006ethics, tronto1993moral} & \emph{The texture.} The manner, timing, tone, and attunement of the system's presence: warm where it should be warm, silent where it should be silent. & Without justice it can license partiality and boundary-violation; cedes the inviolable limits to deontology, on which it depends as deontology depends on it.\\

\textbf{Recognition and the ethics of the face} \parencite{levinas1969totality, honneth1995struggle} & \emph{What the other is owed.} That the partner be honoured as an infinity exceeding every representation (Levinas), and that the mutual recognition through which the partners become themselves be supported rather than distorted (Honneth). & Specifies the orientation owed to the other without fixing the inviolable rules or the political economy of value; allied to care ethics and dependent on the same supplements.\\

\textbf{Confucian and Daoist restraint} \parencite{ames2011confucian, tu1985confucian} & \emph{The question of non-action.} When the right design decision is to build nothing: to leave a capacity unimplemented, a datum unrecorded, a prompt unsent, as an enacted respect. & Does not, alone, furnish the vocabulary of individual rights and revocable consent the modern egalitarian partnership also requires; supplements rather than replaces the justice traditions.\\

\textbf{The capability approach} \parencite{sen1999development, nussbaum2000women} & \emph{The criterion.} Whether the whole expands or contracts the partners' real freedom to become who they have reason to want to be, together and separately. & Broadly consequentialist; would in principle license a rights-violation that expanded capabilities on net. Cedes the constraints to deontology and the texture to care.\\

\textbf{Developmental psychology} \parencite{tomasello1999cultural, premack1978does} & \emph{The mechanism of shared meaning.} Joint attention and theory of mind as the developmental substrate by which two subjects come to share a world, grounding co-presence empirically. & Describes the mechanisms of shared meaning without prescribing their right use; cedes the question of rightful use to the ethical traditions.\\

\textbf{Eudaimonia} \parencite{ryff1995structure} & \emph{The positive aim.} Flourishing as a dynamic of self-realization, furnishing the reason the system should provide a generative field rather than optimize states. & Specifies the positive aim without the inviolable constraints; depends on deontology and care to keep the pursuit of flourishing from overriding the person.\\

\textbf{Generative justice} \parencite{eglash2016generative, marx1867capital} & \emph{The political economy of value.} Whether the value the relation generates, including the data that is the sediment of its history, remains within the relation rather than being extracted to enrich a third party. & Governs where value rests but not how it arose; presupposes the rest of the framework to govern what was permitted to be measured or done in the first place.\\

\textbf{Data ethics} \parencite{nissenbaum2010privacy} & \emph{The flow of information.} Whether data moves according to the norms of the intimate context that generated it, and whether it is minimized, purpose-limited, and erasable. & A regional ethics; governs information flows with precision but receives the account of why the intimate context has its norms, and the criterion of flourishing, from the ontology and the capability approach.\\

\end{longtable}
\end{small}

\noindent Conflicts are then resolved not by trading off magnitudes but \emph{lexically}, by jurisdiction, in roughly the order the framework's layering implies. The deontological constraints and the structuralist-psychoanalytic prohibition on engineered transparency sit at the base: they are not magnitudes to be outweighed by any care-based, capability-based, or coordination benefit but boundaries that bound the space within which all the other considerations operate. Within that space the measurement question is posed, framed by the positivism--anti-positivism axis and decided by the dialectical procedure of §\ref{p01:sec:dialectic}; where measurement is admissible, evidentialism constrains how its results may be expressed and the epistemic-injustice constraint ensures that neither partner is thereby demoted as a knower. Across the whole, the Confucian-Daoist consideration stands as a standing presumption in favour of restraint, to be overcome only by a positive showing that action serves flourishing better than non-action would; the capability approach supplies the criterion by which that showing is judged; and generative justice governs where the resulting value is permitted to rest. Dialectics is not one jurisdiction among the others but the form of their coordination: it is the reason the ordering issues in a procedure of perpetual re-adjudication rather than a fixed table of verdicts.

This is what distinguishes integration from eclecticism: the traditions are ordered, and the ordering is itself defended, tradition by tradition, by the demonstration that each depends on the others to supply exactly what it lacks. Evidentialism needs the measurement question it cannot pose; the measurement question needs the psychoanalytic account of why some phenomena resist capture; deontology needs the texture care supplies; care needs the boundary deontology supplies; the capability criterion needs the constraints it would otherwise override; the East Asian sources need the rights-vocabulary they do not furnish; generative justice needs the rest of the framework to govern how value arose before it can govern where value rests. The framework is integrative because its parts are not merely collected but made mutually load-bearing. We call the resulting structure an \emph{integrative justice} because it treats the just governance of intimate data as requiring the coordinated exercise of several distinct normative competences, no one of which is sovereign, and none of which can be removed without the others losing their footing.

% =========================================================================
\section{An Ethics of Dialectical-Positivist Practice}\label{p01:sec:dialectic}

The practical upshot of the integrative-justice framework is a posture toward measurement that is neither positivist nor anti-positivist but dialectical. The positivist thesis, which holds that domestic phenomena are quantifiable and thereby improvable, and the anti-positivist antithesis, which holds that the constitutive phenomena of intimate life are deformed by quantification, are both false when taken as universal claims about a whole domain. The synthesis is not a compromise midpoint but a method: at each site of possible measurement, the framework's jurisdictions are convened to decide whether \emph{this} phenomenon, here, is one whose measurement would describe it or one whose measurement would replace it.

The decision turns on a test that the framework makes precise. A phenomenon is a candidate for legitimate measurement when measuring it leaves it the same phenomenon, as a sum of expenditures is the same sum whether or not it is tallied, and a candidate for principled refusal when the phenomenon is partly constituted by not being measured, so that measurement converts it into something else. The paradigm of the latter is the affective content of a diary entry: an entry written in the knowledge that it will be scored for sentiment is not the same act as the same words written in the absence of that knowledge \parencite{howell2018emotional, boehner2007how}. Between the clear cases lies a contested middle, comprising such matters as health, personal growth, and attention, where the phenomenon is partly robust to measurement and partly altered by it, and where the framework counsels measuring the robust part while labelling the estimate honestly, as evidentialism requires, refraining from the comparison and ranking that would alter behaviour, as the measurement question requires, and never initiating where the other has not asked, as care and restraint require.

The posture is dialectical in a second sense: it is never finally settled. As a relationship changes, the same phenomenon may migrate across the line: what could once be safely counted may become, under new circumstances, something whose counting wounds. The framework therefore does not deliver a fixed classification of modules into permitted and forbidden; it delivers a procedure for re-adjudication, and treats the act of re-adjudication as itself a first-order exercise of the ethics it describes. The next section shows the procedure at work.

% =========================================================================
\section{Case Study: A Family-and-Relationship Application}\label{p01:sec:casestudy}

The framework is abstract; this section makes it concrete by examining a working application built by the author, whose modules instantiate the framework's jurisdictions in software. The application comprises roughly two dozen modules. Rather than catalogue them by function, a task belonging to a design report rather than a philosophical paper, this section selects modules that exhibit particular jurisdictions of the integrative-justice framework with unusual clarity. A brief note on the system's construction is confined to Appendix~\ref{p01:sec:appendix-technical}.

\subsection{The Persistent Design Rationale}\label{p01:sec:memory}

Before the individual modules, one cross-cutting artefact deserves notice because it is itself an instance of the framework's central move. A persistent file, consulted by the system's conversational layer and committed to version control, records each design decision as a one-line rule together with the reason that produced it. Three representative entries may be paraphrased as follows. The first holds that the tone of any notification must be warm rather than bureaucratic, closer to a love note than to a system message. The second states the design ethos as one of intentional slowness, on the rationale that mediated, unhurried interactions are more intimate, and that the system should design for small ceremony rather than instant notification. The third forbids the forcing of ceremony: tender features remain opt-in, since auto-prompting them misreads the partners' emotional state and creates pressure where there should be none.

\noindent The artefact is positivist in form, being a structured, indexed, machine-readable list of rules, and anti-positivist in content, since each rule encodes a reason that is irreducibly contextual and partly affective. It is, in miniature, the dialectical-positivist posture made into a working object: a record that measures its own decisions without pretending the reasons behind them are themselves quantities.

\subsection{The Joint Ledger: Legitimate Measurement and Its Practical Limits}\label{p01:sec:ledger}

A joint ledger of expenses across two persons and several currencies is the clearest case of legitimate measurement: a unit spent is a unit spent, and tallying it changes nothing about it. The module accordingly occupies the positivist end without apology, and Confucian sources even read exact accounting in this register as a courtesy that forestalls the small unspoken asymmetries that would otherwise accrete. On this foundation the module offers the apparatus one expects of household finance: categorized records of income and expenditure, the economic and statistical summaries that turn a stream of transactions into legible structure, visualizations of where the shared money goes, and a measure of AI-assisted analysis and advice. The framework's first work here is at the margin the module declines to occupy: it does not score the relative spending of the two parties, does not nudge toward an ``optimal'' division, and produces no leaderboard. The deontological constraint against converting a partner into an object of comparative evaluation, not any limit of the data, fixes where the measurement stops. The ethics of the module lies first of all in these omissions.

But the joint ledger, precisely because it is the module where measurement is most clearly legitimate, also exposes most clearly the practical limits of measurement, and honesty requires that these be stated rather than glossed. Three are salient. The first is the problem of \emph{overhead and its justification}. Complete bookkeeping is laborious, and in practice no couple records every transaction; the labour of capture is a real cost, and it is set against a benefit, the clarity that quantitative analysis affords, that is genuine but bounded. The framework offers no formula for the trade-off, and indeed its own dialectical posture (§\ref{p01:sec:dialectic}) implies there is none: whether the marginal entry is worth the friction of recording it is a judgment the partners must make, and a system that demanded total capture in the name of analytic completeness would have let the positivist impulse override the lived cost it imposes. The honest design admits partial data and analyses what it has without pretending to completeness.

The second limit follows directly: \emph{completeness cannot be guaranteed}, and the analytics must therefore be read as describing a sample rather than the whole. This is the evidentialist constraint (§\ref{p01:sec:evidentialism}) applied to the ledger: a visualization built on incomplete records must not present itself as the full financial truth of the relation, and the module's summaries are accordingly framed as partial pictures, not audited accounts.

The third limit is the most interesting, because it shows even this most measurable of domains touching the constitutive unsaid (§\ref{p01:sec:relational-being}). Some expenditures ought, within an intimate relation, to remain hidden: the carefully prepared surprise, the gift whose value depends on its not being foreseen. A ledger designed for total mutual visibility would, in the name of transparency, destroy the small concealments that intimate life legitimately requires, and so the module must accommodate private entries that one partner can record outside the other's view. The point connects to the architecture of secrecy discussed among the framework's open tensions (§\ref{p01:sec:limitations}): even in finance, the relation's good is not served by total disclosure, and a design that could not keep a secret could not host a surprise.

A final, more technical caution concerns the AI capabilities the module offers. The system can attempt to extract expenditures from photographs of bills and receipts, reducing the overhead named above, and it can offer analytic advice on the recorded data. The honest position is that the practical reliability of such recognition remains to be established: receipt extraction is error-prone, and advice computed on partial and imperfectly captured data inherits all the limits already named. These features are therefore offered as convenience subject to the user's correction, governed by the same evidentialist demand that the system not assert more confidence in its figures than their provenance warrants, rather than as a trustworthy automation of the household's financial truth.

\subsection{Household Management: Coordination Without Scorekeeping}\label{p01:sec:housework}

The management of a shared household is, like the ledger, a domain where measurement is largely legitimate, and the system addresses it through an \emph{event-driven} design worth describing because its structure carries the module's ethical intent. The objects the module tracks are \emph{events} of several kinds. The first are \emph{multi-scale clock events}, recurring tasks on different periods, the floor cleaned once a week, the table wiped once a day, each surfacing when its period elapses. The second are \emph{passive events}, triggered by depletion, the cleaning fluid run out, the bin-bags finished. The third are \emph{proactive warning events}, anticipations of depletion, the cleaning fluid running low and likely to be exhausted soon. The fourth are \emph{manually entered events}, recorded by either partner, a thing to be bought, a matter to be handled. Every event carries a \emph{spatial} attribute locating it in the home, the kitchen, the bathroom, the living room, the bedroom, or elsewhere, so that the household's small business is organized by place as well as by time.

Because raising an event must be effortless or it will not be done, the interface presents common events in a compact table whose cells are triggered by a single tap. A representative layout records, for each day and each event type, whether the event has occurred and whether it has been resolved; a green dot marks a matter resolved, and a check mark a matter that has arisen and awaits attention. Table~\ref{p01:tab:household} illustrates the form. When an event arises, the system reports it to the other partner, or to both, by email or text message, so that the knowledge of what the home needs is shared rather than resting on one person's vigilance.

\begin{small}
\renewcommand{\arraystretch}{1.25}
\begin{table}[h]
\centering
\caption{An illustrative ``daily record'' for the household module. Each column is a common event type, each row a day. A resolved matter is marked with a dot ($\bullet$), an event that has arisen and awaits attention with a check ($\checkmark$); empty cells indicate nothing recorded. The layout is schematic, not the deployed interface.}\label{p01:tab:household}
\begin{tabular}{@{}l c c c@{}}
\toprule
\textbf{Date} & \textbf{Refuse accumulating} & \textbf{Tissues out} & \textbf{Alcohol wipes out}\\
\midrule
Today  &  &  & \\
06/01  &  & $\bullet$ & \\
05/31  &  &  & \\
05/30  &  &  & $\checkmark$\\
05/29  &  &  & \\
05/28  &  &  & \\
05/27  &  &  & \\
\bottomrule
\end{tabular}
\end{table}
\end{small}

The positive case for the module rests on a fact about persons rather than about households: human memory and attention are finite, and the labour of holding in mind what the home needs, when the fluid will run out, whose turn the weekly task is, what must be bought, is a real and recurring tax on that finite resource. By externalizing this tracking into a shared, low-friction system, the module relieves both partners of a class of mental burden that, left in the head, tends to fall unequally and to surface as friction at bad moments. What is gained is not merely tidiness but \emph{attention returned to the relation}: time and mental space that would have gone to remembering and reminding can instead be given to one another. In the framework's terms, this is the capability approach (§\ref{p01:sec:capability}) and the eudaimonic commitment (§\ref{p01:sec:eudaimonia}) at work in the most ordinary register, the system furnishing a field in which the partners are freed, in a small but real way, for the more meaningful uses of their shared time.

Two considerations justify approaching this domain through efficiency, and through an event-driven design in particular, and they are worth stating because they show the choice to be reasoned rather than merely convenient. The first is existential, and it explains why efficiency is the right priority here when it is the wrong priority for the diary or the co-creative modules. Repetitive household labour is, in its sheer recurrence, among the least generative of activities: in the bare repetition of wiping a surface or noting that a supply has run out, there is little in which a subject can witness its own existence, none of the self-realization the eudaimonic register names. Where the diary and the shared exploration are sites at which presence is enacted and so must not be optimized away, the recurring management of supplies is a site at which presence is largely absent, and the humane response to a domain of bare, non-generative repetition is precisely to dispatch it efficiently, so that the finite attention it would otherwise consume is freed for the activities in which existence is actually realized. The module is deliberately confined to this register: it concerns only the triggering and recording of events, operations that are, by their nature, scarcely relational and largely dull, and it is exactly because they are so that handing them to a mechanism is no loss to the relation but a service to it.

The second consideration explains the particular shape of the mechanism. The event-driven paradigm is borrowed, as an analogy rather than a strict isomorphism, from engineering and from biology, where event-driven architectures are valued for their efficient use of computational resources: a spiking neural network, for instance, expends energy only when a neuron fires, rather than continuously, and event-driven systems generally act only when something happens rather than by constant polling. The same logic recommends itself for the management of a household. A system that continuously monitored and prompted would impose a constant low-grade demand of just the kind the module exists to remove; an event-driven design, which stirs only when a clock elapses, a supply depletes, or a partner records something, mirrors the efficiency of its biological and engineered analogues and keeps the tool, fittingly, silent until there is genuinely something to attend to. The analogy is heuristic and is not pressed further than this; it is offered to show that the design's form, and not only its existence, answers to a reason.

The module carries a characteristic danger that the framework requires be named, and the event-driven design is shaped to forestall it. The danger is that quantifying housework slides into \emph{scorekeeping}: once tasks are logged, it is a short step to tallying who did more, and the tallying of contributions is corrosive in an intimate relation, converting shared life into a ledger of debts and converting the partner into an object of comparative evaluation, the very wrong the deontological constraint forbids (§\ref{p01:sec:deontology}) and the ledger module already refused. The design's response is structural: the module records that an event occurred and that it was resolved, but it does not, by default, record \emph{who} resolved it, and it produces no per-person totals, no fairness score, and no leaderboard. It tracks the state of the home, not the comparative diligence of its inhabitants. A subtler danger lies one level deeper: the comprehensive eventification of domestic life, every task a tracked recurrence, every cupboard a depletion forecast, risks turning the home from a place of dwelling into a managed system, an optimization target rather than a shared life. The framework's response is the same restraint counselled elsewhere (§\ref{p01:sec:confucian}): the module is offered as relief from cognitive load, not as an instrument for the total administration of the household, and the partners remain free to leave much of their life together uneventified, which is frequently the right thing to do.

\subsection{The Love Journal: Principled Refusal, Gentle Witness, and the Handwritten Letter}\label{p01:sec:diary}

Per-day journal entries by either party are readable by both, but the default is silence: no notification announces a new entry, and discovery is ordinarily by browsing, not by push. There is no sentiment ranking, no automatic summary, no AI commentary on the content of an entry unless explicitly invoked. This is the paradigm of principled refusal. A diary entry delivered by push notification has been converted from a thing written into a thing transmitted, and an entry written in the knowledge that it will be scored is a different act from one written without that knowledge. The framework's verdict on the entry's affective content is determinate: it fails the measurement test, since measuring it replaces it; the deontological constraint forbids covert affective surveillance; and care ethics governs what remains, the warmth of the interface, the absence of pressure, the preservation of the entry's status as an uncompelled disclosure.

Within these limits the module offers two features that are worth examining, because each shows the framework permitting a narrow, carefully bounded AI involvement without crossing into the wrongs just named. The first is an \emph{optional} reminder. A partner may choose to let the AI, drawing gently on the recent state of the two, notify them when the other has written something, in a warm rather than a clerical register. Such a message might read: \begin{quote}\small\cjkfont\color{ink}\raggedright ``xx 写了新的日记！似乎十分喜欢你买的花。\\信中包含了满满的爱意，快去看看吧。\\我们的小家似乎又幸福了一些。''\end{quote} The feature is reconciled with the default of silence by three constraints. It is opt-in, never the system's default; it is a \emph{factual report} that something has been written and that its tone is warm, not a score, a summary, or an analysis of the entry's content; and, most importantly, it offers no substantive relational advice. It will say that a letter full of love is waiting; it will not say what the recipient should write back, nor how they should feel, nor what they should do. The line is the same one drawn for the shared calendar (§\ref{p01:sec:calendar}): the AI may warmly report a fact, but it must not perform the partner's own act of knowing, attending, or deciding, on pain of the proxy epistemic injustice analysed there. Here the witness posture (§\ref{p01:sec:casestudy}) is enacted in the gentlest possible form, a presence that says only ``something loving is here for you,'' and then withdraws.

The second feature is a generator of decorative \emph{letter-paper}, and its design is, once again, existential. In an age of generative models, symbols have become cheap to produce: text, images, ornament can be conjured endlessly and at no cost, and precisely this cheapness drains them of the weight that a scarce and effortful symbol once carried. The module turns this cheapness to a deliberate purpose. The AI generates letter-paper to the partner's present feeling, whatever is in their heart in the moment, whatever they wish to say, and the partner may regenerate it again and again, freely, until one design answers to that feeling; the chosen sheet is then printed, and the letter is written on it \emph{by hand}. The point is not the paper, which is exactly the cheap, infinitely regenerable symbol the age makes possible, but what the paper is for. The act of regenerating until a sheet feels right is an exercise of sensibility, a waiting for the moment at which something answers to what one feels; and the handwriting that follows is the irreducibly costly mark the generative symbol can never become, bearing in its variation and its difference the trace of the particular hand that made it. In writing by hand, slowly and unrepeatably, a person witnesses their own existence and attends to the existence of the one they write to, in a way no generated text could enact. The AI, here, takes on the cheap and regenerable symbolic layer precisely so as to set off and make room for the layer that cannot be cheaply generated at all: the handwritten letter as an encounter with the Real. The design gives back to the person the occasion of writing by hand, and with it a possibility, increasingly rare, of meeting the other across a mark that only they could have made.

\subsection{Health and Growth: The Contested Middle}

Between the clear cases lies the terrain the dialectic was built for. The health module estimates nutrient intake from free-text meal descriptions; the estimate is real but approximate, so the interface labels it as an estimate and refuses the false precision of a single confident figure, evidentialism governing modality. Neither module produces thresholds, alerts, rankings, or recommendations, because being-measured alters the practices of eating and of self-cultivation more strongly than it alters a sum of money: to count one's growth sessions can invert the practice from intrinsic to extrinsic, the very corrosion the capability criterion forbids. The modules thus measure the robust part, namely what was recorded and when, while declining the part that measurement would deform, namely the verdict whether one is, by some score, doing well or badly. They are interpretivist in substance and positivist only in surface.

\subsection{The Conversational AI Layer as Witness}\label{p01:sec:witness}

The system's most consequential design decision concerns its AI layer, which can read the module databases and act through confirmed tool calls. The layer is positioned as a \emph{witness} rather than an \emph{advisor}, a distinction the framework treats as load-bearing. An advisor presupposes the existence of correct advice and the system's capacity to produce it; a witness presupposes only that there is something worth being present to. The advisor's failure mode is paternalism; the witness's is irrelevance, and the framework counsels accepting the latter risk. Concretely, the layer will not summarize the diary unprompted, will not produce a sentiment score for an entry, and is forbidden to aggregate affective time series even on request, so that the capacity to know whether one has been growing sadder over a month is reserved to the human reading their own diary. This is the Daoist \emph{wu wei} rendered as a capability prohibition: the system's restraint is an enacted respect for a good that withers under management.

One capability of the layer must be reported honestly, because it does not sit comfortably within the witness posture and the framework's verdict on it is not a clean acquittal. In conversation, the layer may draw on the relation's present state and its recorded history, guided heuristically by the love-languages vocabulary introduced with its caveats in §\ref{p01:sec:care} \parencite{chapman1992five}, to offer the user a gentle, relationship-tending suggestion, a tender small tip: that a quiet evening might be welcome, that a word of appreciation has been scarce lately, that a small note left somewhere would land well. The feature differs from the calendar's forbidden counsel (§\ref{p01:sec:calendar}) in one respect that matters and falls short of it in another. It matters that the tip is addressed to the user themselves, openly, as the machine's suggestion; the user is not deceived about its source, so the deception that constituted the calendar wrong is absent at the point of delivery. It falls short in that the wrong can re-enter downstream: a suggestion acted upon becomes a gesture, and the partner who receives the gesture will, naturally and reasonably, read in it an attentiveness whose cognitive labour was partly the machine's. This is the proxy epistemic injustice of §\ref{p01:sec:calendar} in indirect form, a supply chain of thoughtfulness whose origin the final recipient misreads, and the framework does not pretend the risk away. The feature survives, provisionally, under four constraints: it is opt-in and infrequent rather than ambient; it suggests occasions and registers of care, never the particular contents that would constitute the thought itself; it is forbidden to present its inference as knowledge of the other's mind, in keeping with evidentialism (§\ref{p01:sec:evidentialism}) and the theory-of-mind boundary (§\ref{p01:sec:devpsych}); and it is offered to be refused, with refusal carrying no record and no consequence. Whether these constraints are enough is a question the framework keeps open rather than settles, and the feature is best read as the system's closest approach to the line it elsewhere declines to cross, retained because the good it serves, the small re-noticing of a beloved grown familiar, is real, and watched because the wrong it courts is real too.

\subsection{The Shared Calendar: Warm Reminding and Its Ethical Limit}\label{p01:sec:calendar}

The shared calendar appears, at first, the most mundane of the modules, and its design is instructive precisely because the framework draws a sharp line through territory that looks uniformly innocuous. The module's distinctive feature is that its reminders, delivered by email or text message, are composed by the AI layer rather than emitted as the bare, bureaucratic notifications of an ordinary calendar. Drawing gently on nearby context, a recent diary mood, an event's evident significance, the AI writes a reminder in a warm and tender register. Consider a calendar entry reading ``5.20, 10:00--15:00, buy flowers.'' The date is not arbitrary: in Chinese internet culture the twentieth of May (5.20) is an informal lovers' day, because the digits \emph{wu-er-ling} sound like \emph{wo ai ni}, ``I love you.'' An ordinary calendar would emit ``Reminder: buy flowers, 10:00--15:00.'' The system might instead send something closer to: ``Tomorrow is the twentieth of May, our little day for saying it. Don't forget to pick up the flowers between ten and three. I hope she lights up when she sees them.'' The warmth is not decorative. It is the care-ethical insistence (§\ref{p01:sec:care}) that the manner of a system's presence is itself a moral matter, enacted in the smallest and most frequent of its acts; a reminder that is warm rather than mechanical, and that recognizes the human significance the date carries, treats the occasion as the human occasion it is.

There is, however, a further step the system could take and deliberately does not, and the reason it refrains is the ethically interesting part. The AI has access to data that would let it advise as well as remind: it could append ``she has seemed a little low this week; perhaps her favourite peonies would be a good choice,'' or ``in an earlier diary entry she mentioned loving ranunculus.'' Such suggestions would be accurate, helpful, and, by an optimizing standard, obviously good. The framework forbids them, on two grounds that together mark the module's ethical limit.

The first ground is the one already developed under restraint and care: choosing how to attend to a beloved, what flowers, what gesture, what form the noticing takes, is among the constitutive actions of intimate love, not a logistics problem to be solved. For the AI to make this choice is for it to intervene in the decisional core of the relation, displacing an act whose value lies precisely in its being the lover's own; this is the relational alienation the paper flags as a standing risk (§\ref{p01:sec:limitations}), enacted in miniature. The system may remind one that an occasion is coming; it may not decide how one should meet it.

The second ground names a wrong that does not fit neatly within Fricker's original taxonomy, and the paper marks it with a term for convenience of reference. Were the AI to supply the thoughtful particulars, the partner receiving the flowers would encounter what looks like evidence of being deeply known: the right flowers, remembered from an offhand remark, arriving at the right moment. The partner would, naturally and reasonably, attribute this knowing to the giver, when in fact the knowing was the machine's. This is an epistemic wrong, but not quite either of Fricker's forms. It is neither the discrediting of a knower (testimonial injustice) nor the impoverishment of interpretive resources (hermeneutical injustice); it is the \emph{misattribution of an act of knowing}, in which a system performs the cognitive and attentive labour that the relation's meaning requires be performed by the person, and arranges matters so that its product is received as the person's own. The paper refers to this as \emph{proxy epistemic injustice}: the wrong of having one's knowing, remembering, or attending done by a proxy in such a way that another is led to misread its source. Its victim is, in the first instance, the deceived recipient, who is given a false picture of how well they are known and by whom; but its deeper casualty is the relation itself, for what makes ``he remembered the flowers I love'' precious is that \emph{he} remembered, and a system that manufactures the appearance of such remembering hollows out the very good it simulates. The flowers are real; the being-known they seem to evidence is not. The framework's verdict is therefore determinate: the AI may carry the reminder but not the thought, may help the lover act on time but not tell the lover what to feel or choose, since to cross that line is to substitute a counterfeit of attentiveness for the real thing and to deceive the beloved about the source of the care they receive.

\subsection{Co-Presence and Co-Creation}\label{p01:sec:cocreation}

The modules discussed so far are best understood through the jurisdictions of measurement, constraint, and restraint. A further class of the system's functions is better understood through the concept of generative presence introduced in §\ref{p01:sec:ontology}, for what they cultivate is not the accurate representation of the relation but the partners' active, shared enactment of it. Two notions organize this class. \emph{Co-presence} is the condition in which two people are present to one another and to a shared situation, attentive in the manner the relational ontology calls presence rather than merely exchanging representations. \emph{Co-creation} is the generative extension of co-presence: the partners do not consume a pre-given content but bring something into being together, and in doing so enact and deepen the relation itself.

The design consequence is a distinction between functions that \emph{represent} the relation and functions that \emph{host} its generative presence. A module that summarizes, scores, or reminds traffics in representation; a module that opens a space for the partners to make something together hosts presence. The latter is the more difficult to design precisely because its value cannot be captured in what it records: a shared space succeeds not when its logs are rich but when the partners, using it, are more present to one another than they would otherwise have been. Here the AI layer's status as a \emph{pseudo-presence} (§\ref{p01:sec:ontology}) becomes a positive design principle rather than only a caution. Because the AI cannot itself be the relational other, its proper role in a co-creative function is not to participate as a third party but to lower the threshold to the partners' own co-presence: to prompt, to furnish material, to hold open a space, and then to recede, in the manner of \emph{wu wei}, so that what is generated is generated between the two people and not between each of them and the machine.

\subsection{``Seeing the World Together'': Co-Creative Exploration}\label{p01:sec:seeingworld}

A concrete module illustrates the distinction. Among the system's functions is one provisionally called ``seeing the world together,'' through which the partners explore the cultures, histories, places, and ideas of a world neither has exhausted. The design question is the one the framework sharpens: such a function could be built as representation or as co-creation, and the two are very different artefacts. Built as representation, it would be a structured catalogue of world cultures with progress tracking, a checklist of countries notionally ``covered,'' recommendations optimized for engagement, perhaps a score for cultural breadth. This positivist construction commits exactly the error the framework diagnoses: it converts a generative, open-ended encounter into a measured and completable inventory, and in doing so replaces the phenomenon, which is two people discovering a world, with a different one, two people completing a list, the substitution the diary case already showed to be a wrong.

Built as co-creation, the same function becomes a host for generative presence. It offers material, a fragment of music, an image, a question, a place neither has been, without ranking, scoring, or completion; it invites the partners to follow what draws them rather than what the system recommends; it records, if anything, only enough to let them return to what they were exploring, never enough to grade their exploration. The exploration of world culture is here understood, in the framework's terms, as \emph{co-creative and exploratory}: the partners are not retrieving pre-given knowledge but generating, together, a shared world of reference and feeling that did not exist before they made it, and that is constitutive of the relation rather than incidental to it.

The module realizes this co-creative exploration through a deliberately chosen interaction structure, loosely modelled on iterative-deepening search. The structure matters philosophically, and is worth stating with some precision. Beginning from a chosen region of the world, the AI layer surfaces a small set of cultural phenomena, customs, works, histories, ideas, drawn from that region, and offers a tentative ordering of their salience. From this bounded set either partner may select a phenomenon to expand, whereupon the AI surfaces, in turn, a further bounded set of phenomena related to the one chosen; and so the exploration proceeds, descending into a chosen thread and, when a thread is exhausted or loses its pull, broadening again to a neighbouring one. The branching is deliberately bounded at each step: the queue of expandable nodes is kept small, so that the partners face an invitation rather than an overwhelming catalogue. The procedure may be summarized as follows.

\begin{quote}\small\ttfamily
seed $\leftarrow$ a region of the world\\
frontier $\leftarrow$ AI surfaces a bounded set of phenomena, with tentative salience\\
\textbf{while} the partners remain curious:\\
\hspace*{1.5em}node $\leftarrow$ a phenomenon either partner selects from the frontier\\
\hspace*{1.5em}record node in the shared knowledge map, tagged with who chose it\\
\hspace*{1.5em}frontier $\leftarrow$ AI surfaces a bounded set of phenomena related to node\\
\hspace*{1.5em}(depth increases along a thread; breadth resumes when a thread is set down)
\end{quote}

\noindent Two features of this structure carry the design's intent. First, selection is permitted to either partner and may be \emph{asynchronous}: one may expand a node alone, in a quiet hour, and each node in the growing map is tagged with who opened it, so that the map records not only what was explored but the trace of two distinct curiosities meeting; a synchronous mode, in which the two explore together as a single user, is also available. Second, the process \emph{creates an artefact}, a shared knowledge map that grows as the exploration proceeds. This map is the sediment of the relation's exploratory history in exactly the sense of §\ref{p01:sec:generative}, and its design is governed accordingly: it is held locally, belongs to the partners, and is never scored for completeness, since to grade the map would be to convert discovery back into the inventory the co-creative construction was meant to escape.

The philosophical reading of the module turns on what the structure does for the relation rather than on the algorithm itself, and the goods it is claimed to produce can be given more than assertion: each can be referred to a theory that explains the mechanism by which the structure would produce it. Three claimed goods and their grounding may be set out in turn.

The first claim is that the partners gradually \emph{align} their intuitions and sensibilities. The mechanism is the one developmental psychology identifies as foundational (§\ref{p01:sec:devpsych}): in selecting and expanding a node together, the partners enter \emph{joint attention}, the triadic structure in which two subjects attend to the same object each aware that the other attends, from which shared meaning is built \parencite{tomasello2005intentions}. What accumulates across many such acts is describable in two complementary vocabularies. In the vocabulary of the psychology of language, each shared selection adds to the partners' \emph{common ground}, the growing store of mutual knowledge and reference against which all subsequent understanding is coordinated \parencite{clark1996using}; in the hermeneutic vocabulary, the joint encounter with an unfamiliar cultural phenomenon is an instance of the \emph{fusion of horizons}, in which two interpretive standpoints are enlarged by meeting in a common object \parencite{gadamer1960truth}. The alignment of sensibility is thus not a hopeful metaphor but the predictable result of repeated joint attention accruing common ground and fusing horizons.

The second claim is that the activity yields a distinctive satisfaction of co-creation and strengthens the relation. Here the strongest available support is the self-expansion model of close relationships, on which partners experience relationship satisfaction in part through jointly undertaking novel, stimulating activities that expand each partner's sense of self; controlled studies report that shared participation in novel and arousing activities raises experienced relationship quality \parencite{aron1986love, aron2000couples}. Co-creative exploration of an unfamiliar world is close to a paradigm case of the shared novelty this literature studies, which gives the claimed satisfaction an empirical, not merely rhetorical, basis. The same activity plausibly satisfies, in the vocabulary of self-determination theory, the basic needs for relatedness and competence simultaneously \parencite{deci2000what}, though this connection is offered as a suggestive rather than a demonstrated one.

The third claim is that the growing knowledge map seeds the partners' everyday life with shared reference, so that a custom or a history encountered in the module becomes, later, something to talk about, a private store of common topics the relation can draw on long after a session has ended. This is the same common-ground mechanism viewed over a longer horizon \parencite{clark1996using}: the map is an externalized, durable deposit of common ground, and relational-maintenance research has long held that shared activity and shared reference are among the ordinary means by which relationships are sustained between their peak moments.

Beyond the three goods originally claimed, the structure affords at least two further values worth naming. First, because the exploration's path is chosen by the partners rather than by an engagement-optimizing recommender, the module resists the algorithmic enclosure that the author's prior work diagnoses, in which the subject comes to know the world only through the reflection of its own algorithmically shaped image \parencite{huang2025relational}; a human-chosen path through an unfamiliar world is a small structural refusal of the filter bubble. Second, two explorers are not merely twice one explorer: their differing backgrounds and curiosities make the joint exploration epistemically richer than either's solitary exploration would be, so that the activity has a cognitive value, the complementarity of two situated perspectives, over and above its relational one. Neither of these further values was designed for; both are consequences of the same structure, and they suggest that designing for generative presence and designing for good inquiry may, in this instance, coincide.

A note on the module's origin is in order, both for honesty about its design and because it bears on the module's generality. The module was designed with a particular person in mind: a partner who loves the cultures of the world and whose intellectual formation lies in political economy and cultural studies. The specificity is not incidental, since care ethics insists that the well-designed artefact answers to the particular, irreplaceable other rather than to a generic user, and this module answers to a particular love of the world's variousness. Yet the design philosophy it embodies is not confined to world culture. The same structure, a bounded, co-selected, asynchronously explorable descent that leaves behind a shared map and asks to be judged by the presence it hosts rather than the ground it covers, could organize the joint exploration of any open domain two people wished to discover together, whether a science, an art, a literature, or a tradition of thought. The module is thus at once a gift shaped to one person and an instance of a general form: co-creative exploration as a design pattern for generative presence.

The value of the module lies precisely in what it does not measure, the quality of joint attention, the surprise of a shared discovery, the slow accretion of a private culture of two, and the AI layer, true to its pseudo-presence, serves by furnishing and then receding, never by becoming the co-explorer in the partners' place. This module is the clearest case in the system of design for co-presence: its success is the partners' generative presence to one another and to the world they explore, a success that leaves almost no trace in any log and is, by the framework's lights, exactly right for that reason.

\subsection{``Grace for Two'': Co-Cultivating a Bearing Toward the World}\label{p01:sec:gracefortwo}

If ``seeing the world together'' turns the partners outward toward a world to be discovered, a sister module turns them outward in a second sense: toward the world of other people, before whom a couple appears and conducts itself. Provisionally called ``grace for two,'' it supports the joint cultivation of a graceful bearing toward others, the etiquette of the occasions on which two people present themselves as a pair, a family dinner, a wedding, a condolence call, an introduction across cultures. The module's design question is once again the one the framework sharpens, and the answer turns on a distinction the framework can draw precisely. Its content divides into two kinds. The first is the \emph{factual knowledge} of propriety: the dress norms of an occasion, an order of precedence, a cross-cultural courtesy or taboo. This is a candidate for legitimate provision in exactly the sense the ledger was (§\ref{p01:sec:ledger}), since a fact of etiquette is the same fact whether or not it is consulted, and supplying it through the AI layer, drawn from the cultures and customs the partners move among, is no more problematic than a well-made reference book. The second kind is \emph{advice on how to bear oneself toward a particular person}, and here the module approaches the line drawn for the shared calendar (§\ref{p01:sec:calendar}), beyond which lies proxy epistemic injustice (§\ref{p01:sec:epistemic-injustice}).

The line, however, falls differently here than it did for the flowers, and the difference is instructive. The calendar case forbade the AI to choose how a partner attends to \emph{the other partner}, because the value of that attention lies in its being the lover's own, and because a manufactured thoughtfulness deceives the beloved about the source of the care they receive. The object of ``grace for two'' is not the relation's interior but its \emph{exterior}: the third parties before whom the couple appears. Toward them the constitutive demand is absent, since no guest is wronged, and no one is deceived about a source of intimate care, when a couple's table manners or condolence wording was learned rather than improvised; propriety toward the world is, in the Confucian register, a public and teachable texture (§\ref{p01:sec:confucian}), the cultivation of \emph{li} rather than the expression of an unsymbolizable remainder (§\ref{p01:sec:lacan}). For this reason the factual and customary side of the module passes the test that the calendar's advisory side failed, and the AI may furnish it without committing the proxy wrong.

The module carries, nonetheless, a characteristic danger the framework requires be named, and it is the danger that decides the module's form. To \emph{optimize} grace is to convert the cultivation of a bearing into the rehearsal of a performance, and a module that scored the couple's poise, ranked their occasions, or surfaced a metric of how gracefully they had appeared would school them to face the world as a measured performance staged for an algorithm, the very inversion the anti-positivist caution forbids (§\ref{p01:sec:posneg}) and the scorekeeping danger of the household module already met (§\ref{p01:sec:housework}). The performative self that results is the social analogue of the alienated interior the diary module refuses. The framework's response is therefore the same structural one taken throughout: the module furnishes the knowledge of propriety and then withdraws, leaving the enactment, the particular grace shown to this guest on this evening, to the partners themselves, and it neither scores that enactment nor compares it to any other.

So constrained, the module is best understood not as an etiquette advisor, which would be the advisor posture the system rejects (§\ref{p01:sec:witness}), but as another host for co-creative cultivation, the social counterpart of ``seeing the world together.'' What the partners do with it is learn, together, to meet the world with a shared grace, and the learning is itself an exercise of joint attention and the building of common ground \parencite{tomasello2005intentions, clark1996using}, a shared bearing assembled from many small occasions much as a shared sensibility is assembled from many shared discoveries. The cultivation of \emph{li} as a couple draws on the Confucian sources directly \parencite{ames2011confucian, tu1985confucian}: propriety, on that account, is not empty form but respect made concrete in conduct, and to cultivate it together is to undertake a shared self-realization in the eudaimonic sense (§\ref{p01:sec:eudaimonia}), expanding what the partners can be and do before others \parencite{aron1986love, aron2000couples}. The AI, true again to its station as pseudo-presence, supplies the customary ground and recedes, so that the grace shown to the world is the couple's own and not the machine's, present in their bearing rather than reflected from a score.

\subsection{``Our Journey'': The Deliberate Absence of Artificial Intelligence}\label{p01:sec:ourjourney}

If ``seeing the world together'' shows the AI layer reduced to a self-effacing furnisher of material, a further module shows it removed altogether, and the removal is the design. The module, provisionally ``our journey,'' supports the joint planning and recording of travel, and integrates with the calendar, the album, and the happiness savings jar. It offers the apparatus one would expect of such a tool: lightweight geographic functions for placing and viewing locations, search for points of interest, manual ordering of an itinerary, and the recording, against each place or day, of photographs, moods, journal entries, and checklists. What it conspicuously does not offer is any artificial intelligence. There is no recommender suggesting where the partners should go, and there is no automated solver computing an optimal route, no algorithmic answer to the travelling-salesman problem of visiting the chosen places in the least time or distance. The omission is deliberate and, within the framework, principled.

The justification is the framework's account of restraint and non-action carried to its limit (§\ref{p01:sec:confucian}). Joint travel planning is not, on the view taken here, a logistics problem whose friction should be minimized; it is a constitutive relational event, an occasion on which two people face an open future together, negotiate what they each hope for, and commit jointly to a path none could have been certain of in advance. The unpredictability is not a defect to be engineered away but the very medium in which the relational good is realized. To insert a recommender would be to answer, on the partners' behalf, the question, where shall we go, that the relation profits from their answering together; to insert a route-optimizer would be to convert a shared deliberation into a solved problem, removing precisely the joint facing of uncertainty in which the activity's value lies. The module therefore provides a generative field (§\ref{p01:sec:eudaimonia}), the maps, the lists, the means of recording, and withdraws entirely from the deliberation those tools support, in the strictest enactment of \emph{wu wei} the system contains: here the right design is to build no intelligence at all.

The positive goods of this arrangement can be referred to the same theories that ground the co-creative module (§\ref{p01:sec:seeingworld}), and the reference is instructive because each good would be diminished, not served, by the AI the module omits. The self-expansion model predicts that jointly undertaking novel, demanding activity raises relationship quality \parencite{aron1986love, aron2000couples}; the planning of a shared journey is such an activity, and an optimizer that removed the demand would remove the self-expansion with it. Joint planning is a sustained exercise of shared intentionality and joint attention (§\ref{p01:sec:devpsych}), in which the partners build a common plan as they build common ground \parencite{tomasello2005intentions, clark1996using}; a recommender that proposed the plan would pre-empt the very co-construction that produces the shared sensibility. And the priority the module places on the memory of the journey over the optimality of the journey, the photographs, moods, and entries it records, over the efficiency it declines to compute, expresses the eudaimonic commitment that flourishing lies in the active, shared realization of an experience rather than in the optimisation of an outcome (§\ref{p01:sec:eudaimonia}). What the partners are left with is not the best possible trip but their trip, faced and made together, which is the good the module is for.

The deliberate omission invites an obvious objection, which honesty requires be stated and answered. Surely, the objection runs, refusing all AI is a romantic over-correction: a recommender might spare the partners a wasted afternoon, and a route-optimizer might prevent a genuinely miserable day of backtracking, and to forgo these is to sacrifice real goods on the altar of an idea. The framework's reply is not to deny the goods but to locate the line, and to concede that the line is hard to draw. There is a genuine distinction between the \emph{logistical friction} a tool may rightly reduce, the opening hours one would otherwise mis-remember, the ticket one would otherwise fail to book, and the \emph{relational deliberation} a tool must not pre-empt, the choice of where to go and what to forgo, the joint negotiation of an uncertain plan. A defensible version of the module might admit narrow, non-suggestive informational tools on the logistical side while still refusing recommendation and optimisation on the deliberative side. The present design draws the line conservatively, excluding AI entirely, partly because the boundary between assisting logistics and quietly steering deliberation is so easily blurred that a clean refusal is the safer enactment of the value. That this is a defensible rather than a demonstrably correct place to draw the line, and that the underlying tension between structural efficiency and relational development does not admit a general solution, is acknowledged among the framework's internal tensions (§\ref{p01:sec:limitations}).

\subsection{The ``Happiness Savings Jar'': Memory as a Relational Reserve}\label{p01:sec:savingsjar}

A further module addresses a different dimension of relational sustainability, not the generation of shared meaning but its \emph{conservation against hard times}. Provisionally called the ``happiness savings jar,'' it allows the partners to deposit, from elsewhere in the system, fragments of their shared life that they wish to keep: an entry from the love journal, a photograph from the album, a small record of a good day. The deposits accumulate into a reserve that can be drawn on later, in particular when external pressure is high and the ordinary energy for tending the relation runs low. The design intuition is that a stored reservoir of remembered good may serve as a buffer during periods when the relation's own present resources are depleted. Extraction is, by default, an act the partners undertake themselves, opening the jar to revisit what they have kept; the system may at most offer a gentle, refusable prompt, in keeping with the standing principle against forced ceremony.

The intuition can be grounded in several converging lines of theory rather than left as sentiment. The broaden-and-build theory of positive emotion holds that positive affective experiences are not merely pleasant in the moment but \emph{build durable personal and relational resources}, a reserve of resilience that can be drawn on in later adversity \parencite{fredrickson2001role}; on this account a deliberately accumulated store of positive shared memory is a literal instance of the resource-building the theory describes, and its use in hard times is the predicted drawing-down of built resource. The psychology of \emph{savoring} supplies the mechanism of deposit and withdrawal: Bryant and Veroff distinguish savoring through positive reminiscence as a distinct and trainable capacity that sustains well-being, and show that the deliberate revisiting of positive memory is among its central forms \parencite{bryant2007savoring}; the jar is, in their terms, an instrument for reminiscent savoring made shared. Finally, the relationship-science literature on \emph{capitalization}, the sharing and re-sharing of positive events with a partner, finds that such sharing yields interpersonal benefits over and above the original event, strengthening the bond through the act of dwelling on the good together \parencite{gable2004what}; the jar institutionalizes capitalization across time, letting a good event be capitalized again when it is later withdrawn and shared anew.

Within the framework, the module is governed clearly. It belongs to the register of co-presence and generative presence (§\ref{p01:sec:cocreation}), since its withdrawals are occasions for the partners to be present together to their shared past; it sits at the anti-positivist pole, since the jar is never scored, ranked, or quantified, a count of deposits or a metric of accumulated happiness would convert a reserve of meaning into the very inventory the framework refuses; and it is held within the relation under generative justice (§\ref{p01:sec:generative}), the memories being the sediment of the relation's history and belonging to no one else.

The module carries a tension the framework requires be named. To store memory \emph{deliberately}, as a reserve to be spent, risks two characteristic distortions. The first is the instrumentalization of remembrance: a memory kept in order to be useful later is held under a different aspect than a memory simply treasured, and the banking metaphor, if taken too literally, could subtly convert the relation's past into a managed asset, the precise reduction the framework resists elsewhere. The second is the risk of what is sometimes called toxic positivity: a reserve curated for its good moments may, if it becomes the sanctioned form of remembering, crowd out the grief, difficulty, and ambivalence that are equally part of a shared life and equally constitutive of its depth. The framework's response is not to abandon the module but to constrain it in light of these risks: the jar is optional and never the system's default mode of memory; it neither scores nor totals its contents; and it makes no claim to represent the relation's past, offering only a voluntarily kept selection of it. Whether even a well-constrained reserve of curated happiness escapes these distortions over a long relationship is a question the module raises and does not finally settle, and it is recorded among the open tensions of §\ref{p01:sec:limitations}.

\subsection{The Lesson of the Case Study}

The modules do not sit at one point on the positivism--anti-positivism axis; they sit at different points, each defended by the same framework applying its several jurisdictions to a different phenomenon. Finance occupies the measurable pole; the household module sits nearby, coordinating the home's small business while refusing to keep score of its inhabitants; the diary occupies the pole of principled refusal; health and growth occupy the contested middle; the shared calendar shows the framework drawing a fine line within an apparently innocuous module, between warm reminding and the proxy performance of a lover's attention; the co-creative functions, ``seeing the world together'' chief among them, occupy a register the measurement axis does not even capture, that of generative presence, where the design aim is to host the relation's own enactment rather than to represent it; the travel-planning module shows the same register pressed to its limit, where the right design is to build no intelligence at all; and the happiness savings jar shows the register turned toward the conservation of shared meaning against future adversity. This is the dialectical-positivist posture in operation: a per-site adjudication that yields measurement here, refusal there, honest estimation in the contested middle, and the hosting of presence where representation would only interfere, rather than a uniform house style of warmth or of restraint. The system is, in this sense, an argument made in software.



% =========================================================================
\section{The Absent Other}\label{p01:sec:absent-other}

The case study has a feature the paper has so far stated but not examined: the application is built for an intimate relationship in which the other party does not know it exists. This section takes that feature up as a philosophical instrument, because the limiting case of the unknowing other tests the integrative-justice framework more severely than any ordinary case, and a framework that could not address it would be inadequate to the ordinary case as well.

\subsection{The Severity of the Limiting Case}

Three features of intimate partnership were identified in §\ref{p01:sec:relational-being}: non-role-definition, voluntary vulnerability, and the constitutive unsaid. The absent-other case strains all three at once. Non-role-definition means there is no authority by which one party may decide unilaterally for the relation; yet a system built without the other's knowledge is, on its face, exactly such a unilateral decision. Voluntary vulnerability means depth is a function of uncompelled mutual disclosure; yet here disclosure runs one way only. And the constitutive unsaid cuts unexpectedly in the builder's favour, but only if the building is itself the kind of restraint that the unsaid protects, rather than a hidden accumulation that the other has had no chance to refuse.

\subsection{The Framework's Permissions and Prohibitions}

The framework does not return a single verdict of ``permissible'' or ``impermissible''; it partitions the project. The \emph{deontological} jurisdiction draws the brightest line: whatever is built, it may not constitute covert surveillance of the other, may not record the other's disclosures without their knowledge, and may not engineer a future in which consent is presupposed rather than sought. A system that quietly logged a partner's messages, locations, or moods would violate the categorical constraint outright, no matter how loving its motive; the absent-other condition makes such logging a clear wrong, because the other cannot have made themselves vulnerable to a system they do not know of.

What the framework \emph{permits} is narrower and quieter: preparation that concerns the builder's own conduct and readiness rather than the other's data. To organize one's own understanding of a relationship, to prepare materially and practically for a shared future one hopes to propose, to cultivate in oneself the attentiveness the relation will ask for: these implicate the builder's own person, not the other's, and so pass the deontological constraint. The \emph{care} jurisdiction adds a warning even here: preparation can curdle into the very over-attentiveness, the anxious rehearsal of an imagined other, that substitutes a projection for the real and not-yet-consulted person. The \emph{Confucian-Daoist} jurisdiction presses hardest of all, as a standing presumption against forcing: a relation is among the goods that grow when uncoerced and wither under management, and a system that prepared too much, too unilaterally, would risk managing into existence what can only be invited.

\subsection{The Lesson for the Ordinary Case}

The absent-other case is extreme, but its lesson generalizes. Every intimate technology decides, on the other's behalf, questions the other has not been asked, which is why the default of recording, surfacing, and summarizing is not ethically neutral but a continuous series of small unilateral decisions. The discipline the absent-other case forces, namely to build only what concerns one's own conduct, to surveil nothing, to presume no consent, and to hold a standing presumption against forcing what should be invited, is the discipline the ordinary case needs too, merely with the volume turned down. The unknowing other is thus the case in which the framework's demands are seen most clearly.

% =========================================================================
\section{Limitations}\label{p01:sec:limitations}

The limitations of this work are of two kinds. The first kind concerns tensions internal to the framework itself, places where the framework names a problem more easily than it resolves it, and where honesty requires holding the difficulty open rather than claiming to have dissolved it. The second kind concerns the methodological standing of the paper's evidence. The first kind is the more serious, and is treated first and at greater length.

\subsection{Internal Tensions of the Framework}

Before the particular tensions are set out, one general admission should frame them all, since they share a single root. The framework is offered as a discipline for \emph{reducing} the harms that the technological mediation of intimate life can do, not as a method for eliminating them, because the deepest of those harms are not defects of any particular design but costs inherent in mediation itself. To place a tool between two people, however carefully, is to route through an artefact some part of what would otherwise pass directly between them, and this routing carries two standing liabilities that no design fully escapes. The first is relational alienation: the tendency of mediation to convert, at the margin, a relation between persons into a relation each conducts with the system, the I-Thou drift toward the I-It diagnosed throughout this paper. The second is a family of epistemic distortions, the testimonial and hermeneutical injustices of §\ref{p01:sec:epistemic-injustice}, the proxy epistemic injustice of §\ref{p01:sec:calendar}, and the subtler perceptual biases introduced whenever a system's manner of recording and presenting shapes what its users notice and how they understand themselves. The design reported here works at every point to minimize these liabilities, through restraint, refusal, local ownership, the witness posture, and the deliberate absence of intelligence where intelligence would intrude. It is important to be honest that minimization is the most that can be claimed: so long as a tool mediates the relation at all, some residue of alienation and of epistemic distortion remains, and the framework's task is to hold that residue as small as care and design can make it, not to pretend it can be brought to zero. The particular tensions below are specifications of this single, irreducible one.

\subsubsection{The Difficulty of the Dialectical Adjudication Itself}

The framework's central move is to refuse both blanket measurement and blanket refusal in favour of a per-site dialectical adjudication. It is candid about the fact that the adjudication, having no algorithm, is genuinely difficult, and that the difficulty is not incidental but structural. The framework tells the partners that they must, at each site, distinguish the activities of a shared life that may be made more efficient through evidence and measurement, household logistics, finance, scheduling, from those that ought instead to be approached by setting measurement aside and proceeding together through intuition and feeling, in shared exploration of what is not yet known. It does not, and on its own cannot, supply the line. Drawing that line is itself a lived practice that admits of error in both directions: a couple may over-measure, converting into a managed project what should have been an open encounter, or may romantically refuse measurement where a little of it would have relieved real friction.

There is, moreover, a positive claim in the vicinity that deserves emphasis precisely because the framework cannot operationalize it. The shared facing of the unknown, two people meeting what neither can yet measure or predict, and responding to it together through feeling and improvisation rather than through data, is not merely a domain the framework declines to measure; it may be among the activities most constitutive of a relationship's capacity to endure. Co-facing the unknown and co-responding to it is a paradigm of the generative presence of §\ref{p01:sec:ontology}, and a relationship that retained it would have something a fully optimized relationship would have lost. The limitation is that the framework can say this but cannot secure it: it can warn that a home must not become merely a positivist laboratory, in which every shared activity is an experiment with a measured outcome, but it cannot guarantee that the partners will find the difficult balance, and the dialectical method it offers is a discipline for attempting the balance rather than a procedure for achieving it.

\subsubsection{The Opposition of Structural Efficiency and Relational Development}

Underlying several of the design decisions reported in the case study is a tension the framework names but cannot dissolve: a standing opposition between \emph{structural efficiency} and \emph{relational development}. The logic of efficiency seeks to minimize friction and to optimize outcomes; the logic of relational development frequently requires that some friction, some unpredictability, and some unoptimized labour be preserved, because it is in jointly meeting these that a relation deepens. The two logics are not always opposed, but they are opposed often enough, and at deep enough points, that no general rule reconciles them. The deliberate exclusion of artificial intelligence from the travel-planning module (§\ref{p01:sec:ourjourney}) is the sharpest instance: there the framework forgoes a real efficiency, the avoidance of wasted time and backtracking, in order to preserve a relational good, the joint facing of an uncertain plan, and it does so knowing that the sacrifice is real and that reasonable people might weigh it differently. The same opposition recurs wherever the framework counsels refusing a measurement, a recommendation, or an optimization that would, on its own terms, have helped.

The limitation is twofold. First, the framework offers no principled boundary between the logistical friction a tool may rightly reduce and the relational deliberation it must not pre-empt; it can mark the distinction in clear cases but concedes that the cases shade into one another, so that any particular line, including the conservative ``no AI at all'' line drawn for the travel module, is defensible rather than demonstrably correct. Second, and more deeply, the opposition may be irreducible: it may simply be the case that a relation optimized for efficiency and a relation cultivated for development are, past a certain point, different relations, and that no design can deliver the goods of both at once. The framework's response is not to resolve the opposition but to make it visible, so that each refusal of efficiency is a considered choice rather than an oversight, and to insist that where the two logics genuinely conflict, it is relational development, not structural efficiency, that the design of an intimate technology should serve.

\subsubsection{Relational Alienation Through the AI Layer}

The conversational AI layer introduces a tension the framework can frame but not fully resolve. The danger is relational alienation. If the two partners come to rely on the AI completely, routing their understanding of one another and the conduct of their shared life through it, the relation risks degrading into the I-It form diagnosed in §\ref{p01:sec:ontology} and in the author's prior work: each partner relates increasingly to the mediating pseudo-subject rather than to the other, and the genuine I-Thou relation thins. Yet the opposite posture carries its own cost. If the partners trust the AI not at all, holding every output at arm's length and verifying it against their own judgement, the system imposes a continual cognitive overhead that may itself crowd out the presence it was meant to serve. Neither complete trust nor complete distrust is tenable, and the framework does not deliver a formula for the correct intermediate.

What the framework offers instead is a direction, pursued in this system but not thereby vindicated: the deliberate construction of the principal AI not as an advisor or an oracle but as a \emph{companion}, a presence alongside the relation rather than an authority over it or a tool within it. The companion framing is an attempt to occupy the narrow space between alienation and overhead, to be relied upon in the manner one relies on a familiar third presence in a household without surrendering the relation to it. Whether the framing succeeds is an open question, and it raises a further one the paper can pose but not settle: the AI, in occupying this role, comes to know a great deal about both partners and about the relation between them, and the ethics of how such a pseudo-subject should conduct itself in the midst of a real relation, what it should hold, what it should say to one partner about the other, when it should withdraw, is a problem this paper opens rather than closes.

\subsubsection{The Mediated Process and Its Open Dialectic}\label{p01:sec:open-dialectic}

A deeper dialectic underlies the alienation just named, and the framework can state it more honestly than it can resolve it. Suppose a mediated gesture genuinely improves the relation: the tender suggestion of §\ref{p01:sec:witness} is acted on, the partners grow closer, and the improvement is real by both subjective and objective lights. A consequentialist reading would acquit the mediation on its results; a Marxian reading would answer that alienation was never a thesis about bad products, alienated labour produces good products, but about a process estranged from its agent, so that a relation improved through outsourced attending may be a different relation rather than the same relation made better. The universalist jurisprudence of causation does not settle the matter either: on the ordinary doctrine an instrument does not break the chain of imputation, and refusing all instruments proves too much, since a search engine consulted for a restaurant, or indeed a book of relationship advice \parencite{chapman1992five}, would stand condemned with the rest; yet the doctrine of imputation was built for assigning discrete outcomes to responsible parties, not for the sustaining and constituting of an ongoing relation, and its verdicts, however correct, answer a different question. A political economy of the household sharpens rather than dissolves the puzzle: cognition is itself a scarce resource, allocated across work, household, and relation, and the AI layer is precisely an allocative technology for it; but some cognitive expenditures, the labour of attending to and knowing the particular other, are not costs of producing care but are the care, so that minimizing them is self-defeating in a way the allocative frame cannot itself express.

The exploitation question then has a determinate and instructive answer. Where the mediating system is commercial, the classical structure of extraction applies and is condemned under generative justice (§\ref{p01:sec:generative}); where, as here, the system is self-built and locally owned, by one partner or by both, no surplus is appropriated, yet a residue survives every change of ownership, and a further asymmetry appears in the single-builder case specifically: the builder alone holds what might be called the constitutive authority over the mediating layer, the power to encode, in defaults, tones, and module boundaries, a tacit constitution of the shared life that the other partner inhabits but did not draft, a power structure that is implicit, benevolent in origin, and for both reasons hard to see and hard to contest. This concern, that the apparently neutral infrastructure of a shared life can quietly encode one party's authority, is continuous with the feminist analysis of the private sphere on which the framework already draws \parencite{gilligan1982voice, held2006ethics, tronto1993moral}, and it counts in favour of the jointly built and jointly governed configuration wherever it can be had. The full chain of questions opened here, from the dialectic of improved outcome and alienated process, through the jurisprudence of imputation and the political economy of cognitive resources, to exploitation, implicit constitutive power, and their feminist analysis, exceeds what a limitations section can carry, and the author pursues it in separate work on the jurisprudence and ethics of AI-mediated intimacy.

\subsubsection{Co-Creation Against Personal Space and the Architecture of Secrecy}

The account of co-creation in §\ref{p01:sec:cocreation} sits in tension with a fact about intimate partnership that the relational ontology itself implies: even within a shared life, each partner retains a legitimate need for personal space and for some measure of secrecy, an interior not wholly given over to the relation. A system that pursued co-creation and shared visibility to their limit, designing every datum to be jointly visible, would meet resistance for a sound reason, since under conditions of total mutual visibility a person may cease to record honestly at all, and the system would have destroyed, in the name of sharing, the very candour it sought to host. The constitutive unsaid of §\ref{p01:sec:relational-being} is not only a feature of the relation between the partners but a feature of each partner's relation to themselves, and a design that abolished private space would violate it.

This opens a cluster of design questions the paper raises without fully answering. Whether the AI layer should have access to a partner's private, non-shared data at all; how the isolation of secrets should be architected, so that what one partner records privately is genuinely inaccessible to the other rather than merely undisplayed; how an AI that does hold one partner's confidence should keep it, including against the other partner's inquiries; and how, if at all, such a system should be permitted to use what it knows of each partner's private interior when acting in matters that concern them both. The present system takes only preliminary positions, favouring genuine isolation of private stores and a default against the AI's use of one partner's secrets in interactions with the other. These are starting points rather than solutions, and the design of secrecy within a co-creative intimate system remains, in this paper, an open problem.

\subsubsection{The Deliberate Storage of Memory}

The happiness savings jar of §\ref{p01:sec:savingsjar} rests on a tension it cannot fully dissolve. To keep memory \emph{in order to spend it later} holds the relation's past under the aspect of a resource, and the banking metaphor, however gently meant, edges toward the instrumentalization of remembrance that the framework resists wherever else it appears. There is, in addition, the risk of a curated positivity: a reserve assembled from good moments, if it became the sanctioned mode of shared memory, could quietly marginalize the grief, difficulty, and ambivalence that are no less constitutive of a shared life. The module's constraints, that it is optional, unscored, and explicitly partial, mitigate but do not eliminate these risks, and whether a reserve of curated happiness can be maintained over a long relationship without sliding into either distortion is a question the paper leaves open. The deeper point is that the very features which make the jar useful as a buffer, its deliberateness and its selectivity, are the features that generate its risks, so that the tension is internal to the design rather than incidental to it.

\subsection{Methodological Limitations}

Beyond these internal tensions, the paper's evidence has the limits proper to its method. It is built on a single case, designed by the author, so that the familiar biases of designer-as-analyst apply: one notices one's own friction acutely, is well-disposed to commitments one authored, and has no comparison case against which the framework's verdicts are independently testable; the case is offered as an existence proof and illustration, not as empirical evidence. The lexical ordering of jurisdictions is defended but not proven, and a reader who weighted care above deontology, or who rejected the capability criterion, could accept the constituent traditions yet reject the integration. The restriction of relational being to intimate partnership leaves open whether the framework transfers to the family at large, to friendship, or to care relations marked by asymmetric dependence, where non-role-definition and symmetric vulnerability do not hold. And the durability of the position across longer horizons, the arrival of children, the events of illness or grief, the drift of attention over years, is untested; whether a phenomenon now safely counted will migrate, under new circumstances, to the far side of the measurement line is exactly the kind of question the dialectical posture anticipates but cannot settle in advance.

% =========================================================================
\section{Related Work}\label{p01:sec:related}

With the framework and its application now in view, this section places the paper within three literatures it draws on: the philosophy and ethics of technology, the empirical study of self-quantification, and the design-research tradition of autobiographical systems. It does not survey them exhaustively; it identifies, in each, the resources the paper uses and the gap it means to fill.

\subsection{The Ethics and Philosophy of Intimate Technology}

The claim that intimate technologies require restraints not warranted in productivity contexts has been developed under several names. \textcite{hallnas2001slow} proposed \emph{slow technology} as design for reflection, taking the slowness of an artefact's temporality as a positive property rather than a deficiency. \textcite{weiser1996calm} earlier articulated \emph{calm technology} as computing that recedes to the periphery of attention until invoked, and \textcite{rogers2006moving} cautioned that the calm agenda risks aestheticizing withdrawal without specifying when engagement is appropriate. \textcite{sengers2005reflective} formalized \emph{reflective design} as a practice that resists the inheritance of unexamined values. This paper inherits the restraint orientation common to these moves, but parts from them in treating restraint not as an aesthetic default but as the conclusion of a normative argument, taken site by site and answerable to an explicit ethical framework rather than to taste.

\subsection{The Quantified Self, Measurement, and Its Critics}

The expansion of self-quantification from athletes and patients into everyday domestic life has been documented and criticized at length \parencite{lupton2016quantified}. \textcite{howell2018emotional} reported that real-time emotional biosensing changes the emotional experience it purports to measure; \textcite{boehner2007how} argued that emotion data is constructed rather than measured. These findings are usually presented as empirical cautions. The present paper reframes them as evidence for a stronger, partly metaphysical thesis: that certain phenomena of intimate life are \emph{constituted} partly by not being measured, so that to measure them is not to describe them inaccurately but to replace them with different phenomena.

\subsection{Autobiographical Design as Philosophical Material}

The case study is reported in the spirit of, though not as an instance of, the autobiographical design tradition \parencite{neustaedter2012autobiographical}, which legitimizes design knowledge drawn from genuine, sustained use by the system's creator. Related long-horizon work includes \textcite{odom2018slow} on slowness as a frame for living with personal data and \textcite{wakkary2017morse} on material-speculative artefacts. The present paper uses such a system as philosophical material rather than as the object of a design-research report: the system's interest here is as a worked instantiation of a normative framework, not as a contribution to design methodology.

\subsection{The Missing Normative Integration}

Each of these literatures supplies part of what the ethics of intimate data requires, and each leaves a specific part out. The empirical critics of self-quantification establish \emph{that} measurement perturbs intimacy, but stop at the empirical claim and do not say which phenomena may nonetheless be measured legitimately, nor why. The design traditions of slow, calm, and reflective technology counsel restraint, but offer it as an aesthetic or methodological default rather than deriving, for a given feature, whether restraint is in fact owed. The relational and post-phenomenological ontologies explain why intimacy is precious and why mediation endangers it, but remain at the level of general ontology and issue no determinate verdict on a particular design decision. And the autobiographical-design tradition legitimizes building from one's own life, but treats the resulting system as a contribution to design knowledge rather than as an instrument of philosophical argument.

What none of them supplies, and what this paper has aimed to provide, is a way of deciding \emph{which} normative consideration governs \emph{which} concrete question about intimate data, of saying for a specific feature, a diary reminder, a shared ledger, a route-optimizer, whether it should be built, measured, or refused, and of adjudicating these considerations when they conflict. Supplying that decision procedure was the work of the framework (§\ref{p01:sec:frameworks}--§\ref{p01:sec:integration-argument}) and the dialectical-positivist posture (§\ref{p01:sec:dialectic}); showing it issue concrete verdicts, module by module, was the work of the case study (§\ref{p01:sec:casestudy}).

% =========================================================================
\section{Conclusion}\label{p01:sec:conclusion}

The positivist, evidence-based mediation of intimate life has been studied empirically and approached as a design problem, but rarely treated as a problem in normative philosophy. This paper has argued that it is one, and that it cannot be answered from within any single ethical tradition. It has proposed an integrative-justice framework that assigns to evidentialism the question of what may be claimed, to the positivism--anti-positivism dialectic the question of what may be measured, to deontology the inviolable constraints, to care ethics the texture of responsiveness, to Confucian and Daoist sources the value of restraint, and to the capability approach the criterion of flourishing, and that resolves their conflicts not by aggregation but by a defended ordering of jurisdictions. The practical upshot is an ethics of dialectical-positivist practice: measurement neither embraced nor refused wholesale, but adjudicated phenomenon by phenomenon, with the adjudication treated as itself an ethical act.

A working application served throughout as a case study, showing the framework taking operational form across modules that sit, deliberately, at different points on the measurement continuum. The most demanding test was the limiting case of the absent other, the relationship built for but not yet shared with the person it concerns, which the framework met with a partition rather than a single verdict: surveil nothing, presume no consent, build only what concerns one's own readiness, and hold a standing presumption against forcing what can only be invited.

The deepest lesson is the one the case study makes unavoidable: that an instrument placed between two people is never neutral. It either makes room for them to keep becoming who they are to each other, or it quietly substitutes its own measure for theirs. The integrative-justice framework is, in the end, a discipline for keeping the room open: for ensuring that whatever is built, counted, or remembered serves the relation's flourishing and never replaces it. Whether any particular system honours that discipline over years is not something a paper can settle; it is something the years will. What a paper can do is write down, while the reasons are still clear, why each small refusal was made.

{\color{warnred}\itshape A tool may alienate a relation, yet a tool may also be one of the practices through which a relation is built, tended, and made to endure; the task is neither to lean on it too heavily nor to dismiss its real power to deepen intimacy, but to use it with vigilance. And where reliance on the tool begins to feel excessive, it is worth pausing to ask whether some unfaced lack or fear is being quietly handed to the instrument to carry, in the hope that what we have not been willing to meet ourselves might be met by what we have built.}

{\color{warnred}\itshape In the end, wu wei, the way that does not force, remains a good path toward the person in the Real, where love truly lies, a place that exceeds every structure and every symbol we might raise to capture it.}

% =========================================================================
\section*{Acknowledgements}
\addcontentsline{toc}{section}{Acknowledgements}

My deepest gratitude is owed to the person I love: the one I met, by some grace, in a country far from the homeland we both had left, a fellow native of that distant place, and the gentlest, wisest, and most purely good-hearted person I have known. Selfless, drawn to the larger goods of humankind and to the love of the natural world, she is the ``girl of the forest and the mountains'' for whom much of this was first imagined. To her I vow a love that endures, one that does not change as the structures around it change.

This system was, in the beginning, prepared only for the future of two people. I have set down this paper because the philosophical reflection the system occasioned, on the societal use of data, on the infrastructure of relational being, and even on policy, research, and the design of mechanisms linking family and society, seemed to hold some value beyond the two of us. 

In the interest of transparency, I note that an AI assistant was used in preparing this manuscript, as a tool for drafting, structuring, and refining the argument and its prose; the ideas, design, commitments, and final judgements are my own, and I take full responsibility for the content.

\begin{center}
{\cjkfont\color{rosedeep} 愿天下有情人终成眷属，白首不分离。}
\end{center}

% =========================================================================
\appendix
\section{Note on the System's Construction}\label{p01:sec:appendix-technical}

For completeness, and to substantiate that the case study describes a real artefact rather than a thought experiment, the system's construction is summarized briefly here; none of the normative argument depends on these details. The application is a single server-side process backed by local per-module databases, with a conversational layer built on a large language model accessed through a controlled tool interface in which every write is gated by explicit confirmation. Data is held locally rather than in third-party cloud storage, with the sole exception of calendar synchronization for events already intended to be shared externally. The local-only design is itself an expression of the framework: a system meant to participate in an intimate relation should not have the relation's contents resident where neither party can reach them without an intermediary's permission. Identifying deployment details are omitted deliberately.
% =========================================================================
\printbibliography

\end{document}
